\documentclass[11pt,a4paper]{article}
\usepackage[utf8]{inputenc}
\usepackage[T1]{fontenc}
\usepackage[top=2.5cm,bottom=2.5cm,left=2.5cm,right=2.5cm]{geometry}
\usepackage{amsmath,amssymb,amsthm,mathtools}
\usepackage[numbers,sort&compress]{natbib}
\usepackage{hyperref}
\usepackage{enumitem}
\usepackage{booktabs}
\usepackage{tikz}
\usetikzlibrary{arrows.meta,positioning,shapes.geometric,fit,decorations.pathreplacing,calc}
\usepackage{xcolor}
\usepackage{float}
\DeclareMathOperator*{\argmin}{arg\,min}
\DeclareMathOperator*{\argmax}{arg\,max}


\theoremstyle{plain}
\newtheorem{theorem}{Theorem}[section]
\newtheorem{proposition}[theorem]{Proposition}
\newtheorem{lemma}[theorem]{Lemma}
\theoremstyle{definition}
\newtheorem{definition}[theorem]{Definition}
\newtheorem{example}[theorem]{Example}
\theoremstyle{remark}
\newtheorem{remark}[theorem]{Remark}

\definecolor{br1}{RGB}{31,119,180}
\definecolor{br2}{RGB}{255,127,14}
\definecolor{br3}{RGB}{44,160,44}
\definecolor{br4}{RGB}{214,39,40}
\definecolor{br5}{RGB}{148,103,189}
\definecolor{br6}{RGB}{140,86,75}
\hypersetup{colorlinks=true,linkcolor=br1,citecolor=br3,urlcolor=br4}
\begin{document}

\title{\textbf{Stochastic Degradation-Aware Optimal Fleet Management with Application to Electric Vehicles}\\[6pt]
{\normalsize A Comprehensive Reference Document}}
\author{L.\ Fiaschi \\ \small ETH Z\"urich \\ \small \texttt{lfiaschi@ethz.ch}}
\date{May 2026}

\maketitle
\begin{abstract}
This is a project scoping where I collect all the needed information to write an article.
The goal is to develop a general purpose model and its algorithmic solver to address the management of a fleet during a certain horizon.
Each member of the fleet is described as a collection of components.
Each day, a member can be sent to serve a mission (route), to maintenance (partial or full repair of each component), or left idle.
The service of a mission damages the components of the fleet member.
Each component has its own stochastic degradation model, and the actual degradation changes according to the mission served.
We show a concrete application to the management of electric-vehicle fleets, such as buses and cars.
\end{abstract}
\tableofcontents
\newpage

\section{Overview of the problem}
\label{sec:overview}

We consider a fleet of $F$ members, each composed of $L$ components subject to stochastic degradation.
The planning horizon spans $2H$ discrete time steps (days).
At each step, every fleet member is assigned to exactly one activity: serving one of $M$ available missions, undergoing maintenance, or remaining idle.
Serving a mission damages the components of the assigned member according to a stochastic degradation process that depends on both the mission and the component.
On a maintenance day, each component may independently undergo imperfect repair (which partially reduces accumulated damage), full replacement (which resets the component to a like-new state), or no intervention at all.

The degradation state of component $\ell$ of fleet member $i$ at time $k$ is described by a random variable $D_{i\ell k}$, whose distribution belongs to a parametric family specified per component.
A component is considered to have failed when $D_{i\ell k}$ exceeds a threshold $\tau$; the schedule must guarantee that the probability of this event remains below $\varepsilon$ at every time step.
To ensure that the schedule can be repeated indefinitely, a \emph{loop constraint} requires that the degradation state at the end of the horizon be stochastically no worse than the state at the mid-point---its probability of failure may not exceed the mid-point one: the second half of the schedule must leave the fleet in at least as good a condition as the first half did.

Table~\ref{tab:notation} collects the notation used throughout this document.
All variables are general and independent of any specific degradation model.

\begin{table}[htbp]
\centering
\caption{Notation and variable definitions.}
\label{tab:notation}
\small
\begin{tabular}{@{}lp{10.5cm}@{}}
\toprule
\textbf{Symbol} & \textbf{Description} \\
\midrule
\multicolumn{2}{@{}l}{\textit{Index sets}} \\
$i \in \{1,\ldots,F\}$        & Fleet member \\
$j \in \{0,1,\ldots,M\}$      & Activity ($j=0$: depot day, i.e.\ maintenance or idle; $j \ge 1$: mission) \\
$k \in \{1,\ldots,2H\}$       & Time step (day) \\
$\ell \in \{1,\ldots,L\}$     & Component \\
\midrule
\multicolumn{2}{@{}l}{\textit{Scalar and fleet-level parameters}} \\
$F$             & Number of fleet members \\
$H$             & Half-horizon length (the full planning window spans $2H$ days) \\
$M$             & Number of missions (must satisfy $F > M$) \\
$L$             & Number of components per fleet member \\
$\tau$          & Degradation failure threshold, taken uniform across components ($\tau = 1$ in the Palmgren--Miner normalisation) \\
$\varepsilon$   & Reliability threshold: $\Pr(D_{i\ell k} > \tau) \le \varepsilon$ for all $i,\ell,k$ \\
\midrule
\multicolumn{2}{@{}l}{\textit{Cost functions}} \\
$C_M$           & Maintenance cost: $\{y_{ik}\}_{i,k} \to \mathbb{R}_{\ge 0}$ \\
$C_R$           & Repair cost: $\{z_{i\ell k}\}_{i,\ell,k} \to \mathbb{R}_{\ge 0}$ \\
$C_{\mathrm{rep}}$ & Replacement cost: $\{x^r_{i\ell k}\}_{i,\ell,k} \to \mathbb{R}_{\ge 0}$ \\
$C_D$           & Damage regularisation cost: $u \to \mathbb{R}_{\ge 0}$ \\
\midrule
\multicolumn{2}{@{}l}{\textit{Component-level parameters}} \\
$\rho_{i\ell}$      & Repair efficiency $\in (0,1]$ \\
\midrule
\multicolumn{2}{@{}l}{\textit{Initial and reference values}} \\
$D^{\mathrm{start}}_{i\ell}$       & Initial degradation state of component $\ell$ of member $i$ (known distribution; must satisfy $\Pr(D^{\mathrm{start}}_{i\ell} > \tau) \le \varepsilon$) \\
$D^{\mathrm{new}}_{i\ell}$ & Post-replacement degradation state (known distribution; typically near-zero damage) \\
\midrule
\multicolumn{2}{@{}l}{\textit{Stochastic state}} \\
$D_{i\ell k}$  & Random variable: degradation of component $\ell$ of member $i$ at time $k$ \\
$\Delta D_{i\ell k}^{(j)}$ & Random damage increment to component $\ell$ of member $i$ when mission $j$ is served on day $k$ \\
\midrule
\multicolumn{2}{@{}l}{\textit{Decision variables}} \\
$x_{ijk} \in \{0,1\}$         & $1$ if member $i$ is assigned to activity $j$ on day $k$ \\
$y_{ik} \in \{0,1\}$          & $1$ if member $i$ spends day $k$ in maintenance (a depot day with at least one intervention) \\
$x^m_{i\ell k} \in \{0,1\}$   & $1$ if component $\ell$ of member $i$ undergoes imperfect repair on day $k$ \\
$x^r_{i\ell k} \in \{0,1\}$   & $1$ if component $\ell$ of member $i$ is replaced on day $k$ \\
$z_{i\ell k} \ge 0$            & Degradation removed by imperfect repair (in expected-damage units) \\
$u \ge 0$                      & Damage regularisation variable \\
\bottomrule
\end{tabular}
\end{table}

\subsection{Scheduling and maintenance variables}

The schedule is encoded by three families of binary decision variables.
The \emph{assignment variable} $x_{ijk}$ determines which activity fleet member~$i$ performs on day~$k$: either serving mission~$j \ge 1$ or spending the day in the depot ($j=0$), either maintenance or idle.
Each member performs at most one activity per day,
\begin{equation}\label{eq:one_activity}
  \sum_{j=0}^{M} x_{ijk} \;=\; 1 \qquad \forall\, i,\,k,
\end{equation}
and each mission must be covered by exactly one member,
\begin{equation}\label{eq:mission_coverage}
  \sum_{i=1}^{F} x_{ijk} \;=\; 1 \qquad \forall\, j \ge 1,\,k.
\end{equation}
Since $F > M$ and mission coverage holds with equality, exactly $F - M$ members are in the depot at every step; the depot headcount is fixed by \eqref{eq:one_activity}--\eqref{eq:mission_coverage}, so the cost model must distinguish \emph{maintenance} days (depot days with at least one intervention, flagged by $y_{ik}$) from idle days, which are free.

On a depot day ($x_{i0k}=1$), the intervention on each component is controlled independently by two additional binary variables, and the \emph{maintenance-day indicator} $y_{ik}$ flags whether any intervention takes place.
The \emph{repair variable} $x^m_{i\ell k}$ indicates whether component~$\ell$ undergoes imperfect repair, and the \emph{replacement variable} $x^r_{i\ell k}$ indicates full replacement.
Both require a maintenance day and are mutually exclusive:
\begin{equation}\label{eq:maint_compat}
  x^m_{i\ell k} \;\le\; y_{ik}, \qquad
  x^r_{i\ell k} \;\le\; y_{ik}, \qquad
  x^m_{i\ell k} + x^r_{i\ell k} \;\le\; 1, \qquad
  y_{ik} \;\le\; x_{i0k}
  \qquad \forall\, i,\,\ell,\,k.
\end{equation}
A component on a depot day with $x^m_{i\ell k} = x^r_{i\ell k} = 0$ receives no intervention (stays idle); its degradation state carries over unchanged.

\subsection{Damage dynamics}

The degradation of each component evolves in discrete time according to the activity assigned to the fleet member that carries it.
We write $\Delta D_{i\ell k}^{(j)}$ for the \emph{random damage increment} sustained by component~$\ell$ of member~$i$ when mission~$j$ is served on day~$k$.
The distribution of each increment is assumed known (from the component's degradation model) and may depend on~$i$,~$\ell$,~$j$, and~$k$.
On days when the fleet member is idle or undergoing maintenance, no damage is accumulated.

At each time step, exactly one of the following mutually exclusive cases applies to every component $(i,\ell)$:

\begin{enumerate}[label=(\roman*)]
  \item \textbf{Mission.} Member~$i$ serves mission~$j \ge 1$ ($x_{ijk}=1$).
    The degradation grows by a mission- and component-dependent random increment:
    \begin{equation}\label{eq:dyn_mission}
      D_{i\ell k} \;=\; D_{i\ell,k-1} \;+\; \Delta D_{i\ell k}^{(j)}.
    \end{equation}

  \item \textbf{Imperfect repair.} The member is on a maintenance day ($y_{ik}=1$) and component~$\ell$ undergoes imperfect repair ($x^m_{i\ell k}=1$).
    The repair reduces the accumulated damage by a fraction determined by the repair efficiency~$\rho_{i\ell}$ and the specific maintenance policy in use.
    Abstractly:
    \begin{equation}\label{eq:dyn_repair}
      D_{i\ell k} \;=\; \mathcal{R}\!\bigl(D_{i\ell,k-1},\;\rho_{i\ell}\bigr),
    \end{equation}
    where $\mathcal{R}$ denotes the repair operator.
    See Section \ref{sec:maint_strategies} for more details.
    The degradation removed by the repair action defines the repair cost variable:
    \begin{equation}\label{eq:z_def}
      z_{i\ell k} \;=\; \mathbb{E}\bigl[D_{i\ell,k-1}\bigr] - \mathbb{E}\bigl[D_{i\ell k}\bigr],
      \qquad \text{when } x^m_{i\ell k}=1.
    \end{equation}
    On days without repair the equality is not imposed and $z_{i\ell k}$ is a free non-negative variable driven to zero by the cost minimisation; equivalently, $z_{i\ell k} = 0$ when $x^m_{i\ell k} = 0$.

  \item \textbf{Full replacement.} The member is on a maintenance day and component~$\ell$ is replaced ($x^r_{i\ell k}=1$).
    The degradation is reset:
    \begin{equation}\label{eq:dyn_replace}
      D_{i\ell k} \;=\; D^{\mathrm{new}}_{i\ell}.
    \end{equation}

  \item \textbf{Idle or maintenance day with no intervention on this component.}
    The degradation is unchanged:
    \begin{equation}\label{eq:dyn_idle}
      D_{i\ell k} \;=\; D_{i\ell,k-1}.
    \end{equation}
\end{enumerate}

Cases~(ii) and~(iii) are mutually exclusive: a component cannot be both repaired and replaced on the same day.
Both require that the member be assigned to a maintenance day ($y_{ik}=1$).
Since exactly one case applies at each step, the dynamics can be written as a single equation:
\begin{equation}\label{eq:dyn_full}
  D_{i\ell k} \;=\;
    D_{i\ell,k-1}\bigl(1 - x^m_{i\ell k}\bigr)\bigl(1 - x^r_{i\ell k}\bigr) + \sum_{j=1}^{M} x_{ijk}\,\Delta D_{i\ell k}^{(j)}
    \;+\; \mathcal{R}\!\bigl(D_{i\ell,k-1},\,\rho_{i\ell}\bigr)\,x^m_{i\ell k}
    \;+\; D^{\mathrm{new}}_{i\ell}\,x^r_{i\ell k},
\end{equation}
with initial condition $D_{i\ell,0} = D^{\mathrm{start}}_{i\ell}$.

\subsection{Stochastic loop constraint}

The planning horizon is split into two halves of length~$H$.
For the schedule to be repeatable indefinitely, the fleet must emerge from the second half in no worse condition than it had at the midpoint.
We formalise this by requiring that the probability of failure at the end of the horizon does not exceed the probability of failure at the midpoint, component by component:
\begin{equation}\label{eq:loop}
  \Pr\!\bigl(D_{i\ell,2H} > \tau\bigr)
  \;\le\;
  \Pr\!\bigl(D_{i\ell,H} > \tau\bigr)
  \qquad \forall\, i,\,\ell.
\end{equation}
Together with the reliability constraint $\Pr(D_{i\ell k} > \tau) \le \varepsilon$ enforced at every step, condition~\eqref{eq:loop} ensures that the second half of the schedule does not consume more reliability margin than the first half created.
A single-threshold comparison of two survival functions is, however, a lossy summary of the state, and by itself it does not guarantee that the inequality persists when the schedule is repeated: the non-invariance that Section~\ref{sec:loop_routeI} exhibits for bound-based certificates arises already for the exact probabilities whenever the state distribution carries more than one free parameter---for a Gaussian state, a variance--mean trade-off can satisfy~\eqref{eq:loop} over one period while the failure probability $\Phi\bigl((\mu-\tau)/\sigma\bigr)$ creeps above~$\varepsilon$ at a later one.
Concatenating copies of the schedule is guaranteed to keep the failure probability below~$\varepsilon$ whenever~\eqref{eq:loop} is certified through a statistic that determines the failure probability monotonically and propagates monotonically under the dynamics~\eqref{eq:dyn_full}: the cumulative shape for the shared-rate Gamma model, the cumulative mean for the shared-dispersion inverse Gaussian model, the componentwise moment condition~\eqref{eq:loop_moments}, or the stronger dominance requirement $D_{i\ell,2H} \le_{\mathrm{st}} D_{i\ell,H}$, which chains because the dynamics is a coordinatewise non-decreasing map of the state and independent randomness (Proposition~\ref{prop:st_closure}).

\subsection{Cost function}

The objective is to minimise the total cost of operating the fleet over the planning horizon.
Four terms contribute to the cost:
\begin{align}\label{eq:objective}
\min \quad
  C_M \left(\{y_{ik}\}_{i,k}\right)
  \;+\; C_{R} \left(\{z_{i\ell k}\}_{i,\ell,k}\right)
  \;+\; C_{\mathrm{rep}}\, (\{x^r_{i\ell k}\}_{i,\ell,k}) +\; C_D\,(u).
\end{align}

All four cost functions are assumed to be non-negative and non-decreasing in each of their arguments: scheduling more maintenance days, removing more degradation, replacing more components, or tolerating higher aggregate damage never decreases the cost.

The first term charges for every maintenance day scheduled across the fleet; it penalises excessive maintenance-day usage (idle depot days are free).
The second term accounts for imperfect repair: whenever a component is repaired, the cost grows with the amount of degradation removed~$z_{i\ell k}$.
The third term charges for every full replacement.

The last term is a damage regularisation penalty.
The variable~$u$ captures the worst-case aggregate expected damage across all fleet members and time steps:
\begin{equation}\label{eq:u_def}
  u \;=\; \max_{i,k}\; \phi\!\Bigl(\mathbb{E}[D_{i1k}],\;\ldots,\;\mathbb{E}[D_{iLk}]\Bigr),
\end{equation}
where $\phi\colon \mathbb{R}^L \to \mathbb{R}_{\ge 0}$ is a non-decreasing aggregation function (e.g.\ the sum or the sum of squares of the expected component damages).
The cost controls the trade-off between minimising operational costs and spreading damage evenly across the fleet.

\subsection{Abstract problem formulation}

Collecting the objective~\eqref{eq:objective}, the scheduling constraints~\eqref{eq:one_activity}--\eqref{eq:maint_compat}, the degradation dynamics~\eqref{eq:dyn_full}, the reliability constraint, and the loop constraint~\eqref{eq:loop}, the fleet management problem reads:

\begin{subequations}\label{eq:abstract}
\begin{align}
\min_{x,\,y,\,x^m,\,x^r,\,z,\,u} \quad
  & C_M\!\bigl(\{y_{ik}\}\bigr)
    + C_R\!\bigl(\{z_{i\ell k}\}\bigr)
    + C_{\mathrm{rep}}\!\bigl(\{x^r_{i\ell k}\}\bigr)
    + C_D(u)
    \label{eq:abs_obj} \\[6pt]
%
\text{s.t.} \quad
  & \sum_{j=0}^{M} x_{ijk} = 1
    && \forall\, i,\,k,
    \label{eq:abs_assign} \\
  & \sum_{i=1}^{F} x_{ijk} = 1
    && \forall\, j \ge 1,\,k,
    \label{eq:abs_cover} \\
  & x^m_{i\ell k} + x^r_{i\ell k} \le y_{ik} \le x_{i0k}
    && \forall\, i,\,\ell,\,k,
    \label{eq:abs_maint} \\[6pt]
%
  & D_{i\ell k} = D_{i\ell,k-1}(1-x^m_{i\ell k})(1-x^r_{i\ell k})
    + \sum_{j=1}^{M} x_{ijk}\,\Delta D_{i\ell k}^{(j)} + \notag \\
  & \qquad\qquad + \mathcal{R}(D_{i\ell,k-1},\rho_{i\ell})\,x^m_{i\ell k}
    + D^{\mathrm{new}}_{i\ell}\,x^r_{i\ell k}
    && \forall\, i,\,\ell,\,k,
    \label{eq:abs_dyn} \\
  & D_{i\ell,0} = D^{\mathrm{start}}_{i\ell}
    && \forall\, i,\,\ell,
    \label{eq:abs_init} \\[6pt]
%
  & \Pr(D_{i\ell k} > \tau) \le \varepsilon
    && \forall\, i,\,\ell,\,k,
    \label{eq:abs_rel} \\
  & \Pr(D_{i\ell,2H} > \tau) \le \Pr(D_{i\ell,H} > \tau)
    && \forall\, i,\,\ell,
    \label{eq:abs_loop} \\[6pt]
%
  & u = \max_{i,k}\;\phi\!\bigl(\mathbb{E}[D_{i1k}],\ldots,\mathbb{E}[D_{iLk}]\bigr),
    \label{eq:abs_u} \\
  & x_{ijk},\, y_{ik},\, x^m_{i\ell k},\, x^r_{i\ell k} \in \{0,1\},\quad
    z_{i\ell k} \ge 0,\quad u \ge 0.
    \label{eq:abs_domain}
\end{align}
\end{subequations}

Problem~\eqref{eq:abstract} is a stochastic combinatorial optimisation problem.
It makes no commitment to a specific degradation model (the distribution of~$\Delta D_{i\ell k}^{(j)}$), repair policy (the operator~$\mathcal{R}$), aggregation function~($\phi$), or cost structure (the functions $C_M, C_R, C_{\mathrm{rep}}, C_D$).
All of these are instantiated when the formulation is specialised to a concrete fleet (Section~\ref{sec:evs}).

\section{Literature review}

\subsection{Stochastic degradation and fatigue models}

Stochastic models of component degradation fall into two broad families.
\emph{Damage-accumulation models} track a cumulative damage variable~$X(t)$ whose distribution evolves over time; failure occurs when $X(t)$ exceeds a threshold.
\emph{Remaining-life models} characterise the number of load cycles a component can sustain before failure, relating applied stress to fatigue life through empirical laws.

\subsubsection{Damage-accumulation processes}

Table~\ref{tab:deg_models} summarises four widely used stochastic processes for modelling cumulative damage.
In each case $X(t)$ denotes the cumulative damage at time~$t$ with $X(0)=0$; all processes have independent, stationary increments (L\'evy processes).

\begin{table}[htbp]
\centering
\caption{Stochastic damage-accumulation models.
Parameters: $\alpha,\beta>0$ (Gamma shape rate / inverse scale), $\mu$ (drift), $\sigma^2$ (diffusion coefficient), $\lambda$ (IG shape rate), $\nu$ (Poisson intensity), $F_Y$ (jump-size distribution).}
\label{tab:deg_models}
\small
\begin{tabular}{@{}p{3cm}p{3.6cm}p{2.2cm}cp{3.5cm}@{}}
\toprule
\textbf{Model} & \textbf{Increment} $X(t)-X(s)$ & \textbf{Marginal} $X(t)$ & \textbf{Mon.} & \textbf{Closure under $+$} \\
\midrule
Gamma process \cite{AbdelHameed1975}
  & $\mathrm{Ga}\!\bigl(\alpha(t{-}s),\,\beta\bigr)$
  & $\mathrm{Ga}(\alpha t,\,\beta)$
  & Yes
  & Same rate~$\beta$ \\[4pt]
Wiener process \cite{WhitmoreSchenkelberg1997}
  & $\mathcal{N}\!\bigl(\mu(t{-}s),\,\sigma^2(t{-}s)\bigr)$
  & $\mathcal{N}(\mu t,\,\sigma^2 t)$
  & No
  & Always \\[4pt]
Inverse Gaussian process \cite{YeChen2014}
  & $\mathrm{IG}\!\bigl(\mu(t{-}s),\,\lambda(t{-}s)^2\bigr)$
  & $\mathrm{IG}(\mu t,\,\lambda t^2)$
  & Yes
  & Same $\mu^2/\lambda$ \\[4pt]
Compound Poisson \cite{EsaryMarshallProschan1973,Wang2022}
  & $\sum_{n=N(s)+1}^{N(t)} Y_n$
  & $\mathrm{CP}(\nu t,\,F_Y)$
  & If $Y{\ge}0$
  & Always \\
\bottomrule
\end{tabular}
\end{table}

\paragraph{Gamma process.}
Introduced by Abdel-Hameed~\cite{AbdelHameed1975} and surveyed in depth by van~Noortwijk~\cite{vanNoortwijk2009}, the Gamma process is the most widely used model for monotonic degradation (wear, corrosion, crack growth).
The process $\{X(t)\}$ has non-negative independent increments $X(t)-X(s) \sim \mathrm{Ga}(\alpha(t-s),\,\beta)$, where $\alpha>0$ is the shape rate and $\beta>0$ is the inverse scale (rate).
Its mean and variance at time~$t$ are $\mathbb{E}[X(t)] = \alpha t/\beta$ and $\mathrm{Var}[X(t)] = \alpha t/\beta^2$.
The Gamma family is closed under convolution when the rate parameter~$\beta$ is shared: if $X_1 \sim \mathrm{Ga}(\alpha_1,\beta)$ and $X_2 \sim \mathrm{Ga}(\alpha_2,\beta)$ are independent, then $X_1+X_2 \sim \mathrm{Ga}(\alpha_1+\alpha_2,\beta)$.

\paragraph{Wiener process (Brownian motion with drift).}
The Wiener degradation model~\cite{WhitmoreSchenkelberg1997} sets $X(t) = \mu t + \sigma W(t)$, where $W(t)$ is a standard Brownian motion, $\mu>0$ the drift rate, and $\sigma>0$ the diffusion coefficient.
Increments are normally distributed: $X(t)-X(s) \sim \mathcal{N}(\mu(t-s),\,\sigma^2(t-s))$.
The normal family is unconditionally closed under addition.
Unlike the other three models, the Wiener process is \emph{not} monotone: sample paths can decrease, making it suitable for degradation phenomena with short-term fluctuations.
A notable property is that the first-passage time to a fixed threshold~$\tau$ follows an Inverse Gaussian distribution: $T_\tau \sim \mathrm{IG}(\tau/\mu,\;\tau^2/\sigma^2)$.

\paragraph{Inverse Gaussian process.}
Proposed as a degradation model by Ye and Chen~\cite{YeChen2014}, the IG process has non-negative increments $X(t)-X(s) \sim \mathrm{IG}(\mu(t-s),\,\lambda(t-s)^2)$, with mean $\mu t$ and variance $\mu^3 t/\lambda$.
Its increments are heavier-tailed than those of the Gamma process, making it appropriate when occasional large damage jumps are observed.
The IG distribution is closed under convolution when the dispersion parameter $\mu^2/\lambda$ is shared across summands; however, the sum of two independent IG \emph{processes} with different parameters is not itself an IG process~\cite{YeChen2014}.

\paragraph{Compound Poisson process.}
The compound Poisson process $X(t) = \sum_{n=1}^{N(t)} Y_n$ models shock-based degradation~\cite{EsaryMarshallProschan1973}, where shocks arrive according to a Poisson process $N(t)$ with rate~$\nu$ and each shock causes a random damage~$Y_n$ drawn i.i.d.\ from a distribution~$F_Y$.
When $Y_n \ge 0$, the process is monotone.
The sum of independent compound Poisson processes is always a compound Poisson process.
This model is natural for discrete damage events (e.g.\ impacts, overloads) rather than continuous wear; see Wang~\cite{Wang2022} for a recent application to structural reliability.

\subsubsection{Discrete-time formulation}
\label{sec:discrete_deg}

The continuous-time processes of Table~\ref{tab:deg_models} are sampled at discrete steps in the fleet management setting of Section~\ref{sec:overview}: each day $k$, the damage of component~$(i,\ell)$ grows by a random increment $\Delta D_{i\ell k}^{(j)}$ whose distribution depends on the mission~$j$ served that day.
The cumulative damage after $k$~steps is therefore a random walk
\begin{equation}\label{eq:discrete_damage}
  D_{i\ell k} \;=\; D_{i\ell,0} \;+\; \sum_{t=1}^{k} \Delta D_{i\ell t},
\end{equation}
where the increments $\Delta D_{i\ell t}$ are independent but \emph{not} identically distributed (different missions may be served on different days).
Table~\ref{tab:discrete_models} restates the four models in this discrete-time setting.

\begin{table}[htbp]
\centering
\caption{Discrete-time specialisation of the damage-accumulation models.
Each row gives the distribution of a single-step increment~$\Delta D$ and the resulting marginal of $D_k = \sum_{t=1}^{k}\Delta D_t$ when all increments share the parameters required for closure.}
\label{tab:discrete_models}
\small
\begin{tabular}{@{}p{2.1cm}p{3.4cm}p{4.0cm}p{4.2cm}@{}}
\toprule
\textbf{Model} & \textbf{Increment} $\Delta D_t$ & \textbf{Marginal} $D_k$ (homogeneous) & \textbf{Closure condition} \\
\midrule
Gamma
  & $\mathrm{Ga}(\alpha_t,\,\beta)$
  & $\mathrm{Ga}\!\bigl(\sum_{t}\alpha_t,\;\beta\bigr)$
  & Same rate~$\beta$ \\[4pt]
Gaussian
  & $\mathcal{N}(\mu_t,\,\sigma_t^2)$
  & $\mathcal{N}\!\bigl(\sum_{t}\mu_t,\;\sum_{t}\sigma_t^2\bigr)$
  & Always \\[4pt]
Inv.\ Gaussian
  & $\mathrm{IG}(\mu_t,\,\lambda_t)$
  & $\mathrm{IG}\!\bigl(\sum_{t}\mu_t,\;\frac{(\sum_t \mu_t)^3}{ \sum_t(\mu_t^3/\lambda_t)}\bigr)$
  & Same $\mu^2/\lambda$ \\[4pt]
Compound Poisson
  & $\sum_{n=1}^{N_t} Y_n$
  & $\mathrm{CP}\!\bigl(\sum_t\nu_t,\,F_Y\bigr)$
  & Same $F_Y$ \\
\bottomrule
\end{tabular}
\end{table}

The closure conditions determine whether the marginal distribution of~$D_{i\ell k}$ remains within the same parametric family after summing increments from different missions.
When the condition holds, the distribution of~$D_{i\ell k}$ is available in closed form as a function of the cumulative parameters, which is essential for evaluating the reliability constraint $\Pr(D_{i\ell k} > \tau) \le \varepsilon$ efficiently.
When it does not hold---for instance, summing Gamma increments with different rate parameters---the marginal must be computed numerically or approximated.

In the fleet management problem of Section~\ref{sec:overview}, each day brings an increment with mission-dependent parameters.
Closure therefore requires that the parameter shared across missions is indeed constant for a given component: the rate~$\beta$ for Gamma, or the dispersion $\mu^2/\lambda$ for Inverse Gaussian.
The shape (Gamma) and mean/variance (Gaussian) parameters may vary freely across missions and accumulate additively; for the Inverse Gaussian the means accumulate additively and the cumulative shape is induced by the shared dispersion, $\lambda = (\sum_t \mu_t)^2\big/(\mu^2/\lambda)$.
For the Gaussian model, no restriction is needed: normal random variables are unconditionally closed under addition.

\subsubsection{Remaining-life models: the S--N curve and Palmgren--Miner rule}

The stress--life (S--N) approach, originating with W\"ohler's experiments in the 1850s, relates the stress amplitude~$\sigma_a$ to the number of cycles to failure~$N_f$.
The most common analytical form is Basquin's law~\cite{Basquin1910}:
\begin{equation}\label{eq:basquin}
  \sigma_a^{\,m}\, N_f \;=\; C,
\end{equation}
where $m>0$ is the fatigue exponent and $C>0$ is a material constant, both typically fitted from experimental data.
Equivalently, $N_f = C\,\sigma_a^{-m}$, which appears as a straight line on a log--log plot.

Under variable-amplitude loading, the Palmgren--Miner linear damage rule~\cite{Palmgren1924,Miner1945} defines the cumulative damage fraction as
\begin{equation}\label{eq:miner}
  D_{\mathrm{PM}} \;=\; \sum_{i} \frac{n_i}{N_{f,i}},
\end{equation}
where $n_i$ is the number of applied cycles at stress level~$\sigma_{a,i}$ and $N_{f,i}$ is the corresponding life from the S--N curve.
Failure is predicted when $D_{\mathrm{PM}} \ge 1$.
The rule assumes that damage accumulates linearly and independently of load ordering---an approximation that has been questioned experimentally (the critical damage sum $D_{\mathrm{PM}}^*$ at failure is scattered around~1 rather than being exactly~1) and extended stochastically by treating $D_{\mathrm{PM}}$ and/or $D_{\mathrm{PM}}^*$ as random variables~\cite{Shen1980,CastilloFernandezCanteli2009,Hilow2017,BohmEtAl2021,MunizCalvente2022}.

Two sources of randomness enter the Palmgren--Miner accumulation:
\begin{enumerate}[label=(\roman*)]
  \item \emph{S--N scatter.}
    At a given stress level~$\sigma_a$, the fatigue life~$N_f$ is not a deterministic constant but a random variable whose distribution reflects material variability, surface finish, and testing conditions.
    The increment $\Delta D_{\mathrm{PM}} = n/N_f$ therefore inherits this randomness.
  \item \emph{Random critical damage sum.}
    Experimental evidence shows that the critical value $D_{\mathrm{PM}}^*$ at failure is itself scattered (typically in the range $0.5$--$2.0$) rather than being exactly~1~\cite{Shen1980}.
    This can be modelled by replacing the deterministic threshold with a random one, or equivalently by normalising the damage by a random capacity.
\end{enumerate}
In the fleet management framework of Section~\ref{sec:overview}, source~(i) is the relevant one: the threshold is fixed at $\tau = 1$ and the randomness enters through the increments.
We therefore focus on the distributional assumptions for~$N_f$ and the resulting distribution of~$\Delta D_{\mathrm{PM}}$.
Table~\ref{tab:sn_scatter} summarises the three most common choices.

\begin{table}[htbp]
\centering
\caption{Distributional assumptions for S--N scatter and their effect on the Palmgren--Miner increment $\Delta D_{\mathrm{PM}} = n/N_f$.
In all cases, $\mu(\sigma_a)$ and the scale/shape parameters depend on the stress level through Basquin's law~\eqref{eq:basquin}.}
\label{tab:sn_scatter}
\small
\begin{tabular}{@{}p{2.1cm}p{3.4cm}p{3.5cm}p{1.2cm}p{3.0cm}@{}}
\toprule
\textbf{Model} & \textbf{Life} $N_f$ & \textbf{Increment} $\Delta D_{\mathrm{PM}}$ & \textbf{Mon.} & \textbf{Closure under $+$} \\
\midrule
Log-normal
  & $\ln N_f \sim \mathcal{N}\!\bigl(\mu,\,\sigma^2\bigr)$
  & $\mathrm{LogN}\!\bigl(\ln n {-} \mu,\,\sigma^2\bigr)$
  & Yes
  & No \\[4pt]
Weibull
  & $N_f \sim \mathrm{Wei}(k,\,\lambda)$
  & Inv.\ Weibull (Fr\'echet)
  & Yes
  & No \\[4pt]
Gumbel
  & $\ln N_f \sim \mathrm{Gumbel}(\mu,\,\beta)$
  & Log-Gumbel
  & Yes
  & No \\
\bottomrule
\end{tabular}
\end{table}

\paragraph{Log-normal.}
The most widely used assumption~\cite{MunizCalvente2022}: at a fixed stress level~$\sigma_a$, the fatigue life satisfies $\ln N_f \sim \mathcal{N}(\mu(\sigma_a),\,\sigma^2)$, where $\mu(\sigma_a) = \ln C - m\ln\sigma_a$ from Basquin's law and $\sigma$ captures the scatter.
Since $\ln(n/N_f) = \ln n - \ln N_f$, each increment is log-normally distributed: $\Delta D_{\mathrm{PM}} \sim \mathrm{LogN}(\ln n - \mu,\,\sigma^2)$, with mean $n\,e^{-\mu + \sigma^2/2}$ and variance $n^2 e^{-2\mu+\sigma^2}(e^{\sigma^2}-1)$.
The log-normal family is \emph{not} closed under summation: the sum of independent log-normals has no standard parametric form.
In practice, $D_{\mathrm{PM},k} = \sum_{t=1}^{k}\Delta D_{\mathrm{PM},t}$ is often approximated as log-normal via moment-matching (Fenton--Wilkinson method), but this is an approximation whose quality degrades as the number of summands grows or the coefficient of variation increases.

\paragraph{Weibull.}
Widely adopted in fatigue standards (e.g.\ Eurocode, ASTM): $N_f \sim \mathrm{Wei}(k,\,\lambda(\sigma_a))$ with shape~$k$ and scale~$\lambda$ depending on the stress level.
The reciprocal $1/N_f$ follows an inverse Weibull (Fr\'echet) distribution, so $\Delta D_{\mathrm{PM}} = n/N_f$ is a scaled Fr\'echet variable.
The sum of independent Fr\'echet variables has no closed form, making the marginal of~$D_{\mathrm{PM},k}$ intractable analytically.

\paragraph{Gumbel (extreme-value type~I).}
Castillo and Fern\'andez-Canteli~\cite{CastilloFernandezCanteli2009} proposed a compatible statistical model for the entire S--N field based on a Gumbel (minimum) distribution for~$\ln N_f$: $\ln N_f \sim \mathrm{Gumbel}(\mu(\sigma_a),\,\beta)$.
The increment $\Delta D_{\mathrm{PM}} = n/N_f$ then follows a log-Gumbel distribution.
As with the other two families, the log-Gumbel is not closed under summation.

\paragraph{Contrast with damage-accumulation models.}
A fundamental difference with the processes of Table~\ref{tab:deg_models} is that \emph{none} of the standard distributional assumptions for S--N scatter yields increments that are closed under summation.
The marginal distribution of $D_{\mathrm{PM},k}$ after $k$~loading blocks must therefore always be computed via convolution (analytically intractable in general), Monte Carlo simulation, or moment-based approximations.
In the fleet management formulation, this means that the reliability constraint $\Pr(D_{\mathrm{PM},k} \ge 1) \le \varepsilon$ cannot be evaluated in closed form from the cumulative parameters alone, in contrast to the Gamma or Gaussian models where the CDF of the marginal is directly available.
The upper bounds of Section~\ref{sec:stoch_orders} (Markov, Cantelli, Chernoff) remain applicable: because the mean and variance of~$D_{\mathrm{PM},k}$ are additive over increments regardless of the distributional family, the Cantelli bound provides a conservative but tractable constraint.
Alternatively, if the scatter is small enough that a deterministic S--N curve (e.g.\ the mean or a design percentile) is an acceptable approximation, the problem reduces to the deterministic Miner's rule and all constraints become linear.

\subsubsection{Conservative surrogates for the empirical damage-increment distribution}
\label{sec:cons_surrogates}

In the remaining-life setting the per-mission increment is the Palmgren--Miner contribution $\Delta D^{(j)} = \sum_i n_i/N_{f,i}$ of Equation~\eqref{eq:miner}, and in our intended use its distribution is not postulated but \emph{measured}: a rainflow count~\cite{MunizCalvente2022} of load histories collected for mission~$j$ yields an empirical probability density of the non-negative increment $\Delta D^{(j)}$.
We deliberately do \emph{not} fit a parametric law to this density---a fit ties the model to one data set and one functional form, undermining the generality the abstract formulation seeks---and instead summarise it by a few distributional functionals (moments, support, selected quantiles) and replace it by a \emph{conservative} surrogate.
Conservative here means that the surrogate never \emph{under}-states damage, so that the reliability constraint $\Pr(D_{i\ell k} > \tau) \le \varepsilon$ and the loop constraint~\eqref{eq:loop} remain valid by construction.
Two features of the problem delimit the admissible options.
First, both constraints act on the \emph{accumulated} damage $D = \sum_t \Delta D_t$ of Equation~\eqref{eq:discrete_damage}, not on a single increment, so the surrogate must be evaluated after summation.
Second, as observed in the contrast above, none of the standard S--N scatter families is closed under summation, so the law of $D$ is generally unavailable in closed form and must be either bounded or replaced by a convolution-closed surrogate.

Crucially, Route~I never requires the moments or support of the accumulated damage~$D$ in advance: it consumes only \emph{per-mission} information, which the schedule then assembles through the assignment counts.
Fix a component $(i,\ell)$ and an accumulation window, write $D = D_{i\ell}$ for its damage over that window, and let $n_{ij} = \sum_t x_{ijt}$ be the number of days mission~$j$ is served.
Modelling repeated executions of mission~$j$ on component~$\ell$ as independent draws from a mission-specific increment law---with mean $\mu_{\ell j}$, variance $v_{\ell j}$, support $[0,b_{\ell j}]$, and cumulant generating function $\psi_{\ell j}(s) = \ln \mathbb{E}\big[e^{s\Delta D^{(j)}}\big]$, where $\mathbb{E}[e^{sX}]$ is the moment-generating function---independence across days gives
\begin{equation}\label{eq:additive_functionals}
  \mathbb{E}[D] = \sum_j n_{ij}\,\mu_{\ell j}, \qquad
  \operatorname{Var}[D] = \sum_j n_{ij}\,v_{\ell j}, \qquad
  \ln M_D(s) = \sum_j n_{ij}\,\psi_{\ell j}(s).
\end{equation}
Each functional of~$D$ is therefore a count-weighted sum of per-mission quantities---linear in the schedule, with the support entering only through the per-mission $b_{\ell j}$, never through a shared worst-case value.
Repetition \emph{merges} into the tail description additively: serving mission~$j$ a total of $n_{ij}$ times multiplies its cumulant generating function by $n_{ij}$, and distinct missions combine by summation, so a gentle mission keeps its small $\psi_{\ell j}$ beside a punishing one rather than being averaged with it.
These per-mission functionals support two routes.
\emph{Route~I} bounds the upper tail of $D$ directly from the functionals in~\eqref{eq:additive_functionals}.
\emph{Route~II} replaces each increment by a stochastically larger surrogate drawn from a family closed under convolution and propagates the order to the sum, using the following standard closure property of the usual stochastic order $\le_{\mathrm{st}}$ (Definition~\ref{def:st}).

\begin{proposition}[Closure of $\le_{\mathrm{st}}$ under independent convolution]\label{prop:st_closure}
Let $\Delta D_t \le_{\mathrm{st}} \widetilde{\Delta D}_t$ for $t = 1,\dots,k$, with $\{\Delta D_t\}_t$ mutually independent and $\{\widetilde{\Delta D}_t\}_t$ mutually independent.
Then $\sum_t \Delta D_t \le_{\mathrm{st}} \sum_t \widetilde{\Delta D}_t$.
More generally, the conclusion holds with the sum replaced by any coordinatewise non-decreasing map $g$ (see Shaked and Shanthikumar~\cite[\S1.A]{ShakedShanthikumar2007}).
\end{proposition}
\begin{proof}
By the coupling characterisation of $\le_{\mathrm{st}}$ (Definition~\ref{def:st}), for each $t$ there exist $\hat X_t \stackrel{d}{=} \Delta D_t$ and $\hat Y_t \stackrel{d}{=} \widetilde{\Delta D}_t$ on a common probability space with $\hat X_t \le \hat Y_t$ almost surely.
Realise these couplings independently across $t$.
Then $\sum_t \hat X_t \le \sum_t \hat Y_t$ almost surely, while independence gives $\sum_t \hat X_t \stackrel{d}{=} \sum_t \Delta D_t$ and $\sum_t \hat Y_t \stackrel{d}{=} \sum_t \widetilde{\Delta D}_t$; hence $\sum_t \Delta D_t \le_{\mathrm{st}} \sum_t \widetilde{\Delta D}_t$.
Any non-decreasing $g$ preserves the almost-sure inequality, yielding the general statement.
\end{proof}

Proposition~\ref{prop:st_closure} is what makes Route~II usable beyond a single sum: the dynamics~\eqref{eq:dyn_full} is a coordinatewise non-decreasing map of the increments and the (independent) reset variables, so a per-increment stochastic upper bound propagates to a bound on the entire state $D_{i\ell k}$---repairs and replacements included---and hence to $\Pr(D_{i\ell k} > \tau)$ at every step.
Table~\ref{tab:cons_surrogates} collects the resulting options together with their guarantees and trade-offs.

\begin{table}[htbp]
\centering
\caption{Ways to summarise the empirical per-mission damage-increment distribution and propagate it to the accumulated damage $D = \sum_t \Delta D_t$ for the constraint $\Pr(D > \tau) \le \varepsilon$.
``Conservative'' means a guaranteed upper bound on the failure probability (Route~I) or $\le_{\mathrm{st}}$ dominance (Route~II); ``$\widetilde{D}$'' means the method yields an explicit surrogate \emph{law} for the sum---needed, e.g., for the loop constraint, which compares two survival functions---as opposed to a one-sided bound.
All inputs are \emph{per mission}; the schedule assembles them through the assignment counts $n_{ij}$ via Equation~\eqref{eq:additive_functionals}, so the support enters only per mission and is never shared.}
\label{tab:cons_surrogates}
\footnotesize
\setlength{\tabcolsep}{4.5pt}
\begin{tabular}{@{}p{1.7cm}p{1.9cm}p{0.9cm}p{0.7cm}p{2.9cm}p{2.9cm}p{1.5cm}@{}}
\toprule
\textbf{Method} & \textbf{Inputs} & \textbf{Cons.?} & \textbf{$\widetilde{D}$?} & \textbf{Pros} & \textbf{Cons} & \textbf{Ref.} \\
\midrule
\multicolumn{7}{@{}l}{\textit{Route I --- bound the tail of $D$ directly (no surrogate law)}} \\
Markov
  & mean
  & Yes & No
  & needs only the mean; always valid
  & extremely loose
  & \cite{BertsimasPopescu2005} \\[3pt]
Cantelli
  & mean, var
  & Yes & No
  & only two moments; tight \emph{as} a two-moment bound
  & ignores independence; loose in the tail
  & \cite{BertsimasPopescu2005} \\[3pt]
Bernstein
  & per-mission mean, var, support $b_j$
  & Yes & No
  & exploits independence; few parameters; moment inputs linear in the counts
  & needs per-mission support; one-sided only
  & \cite{BoucheronLugosiMassart2013} \\[3pt]
Chernoff
  & per-mission c.g.f.\ $\psi_j(s)$ (or mean $+$ support)
  & Yes & No
  & tightest in the tail; count-weighted; valid for any fixed $s>0$
  & needs the c.g.f.; empirical c.g.f.\ is an estimate, not a bound
  & \cite{Chernoff1952} \\[3pt]
\multicolumn{7}{@{}l}{\textit{Route II --- dominate each increment, then convolve (explicit $\widetilde{D}$)}} \\
Moment-matched Gamma / IG
  & mean, var
  & No & Yes
  & explicit law; closed under $+$; ordering tables apply
  & matches moments but does not dominate (survival curves cross)
  & \cite{vanNoortwijk2009,YeChen2014} \\[3pt]
Dominating shared-rate Gamma
  & moments $+$ empirical survival
  & Yes & Yes
  & conservative \emph{and} explicit; serves reliability and loop; shape linear in the schedule
  & shared rate $+$ dominance margin add slack; commits to Gamma
  & \cite{ShakedShanthikumar2007} \\
\bottomrule
\end{tabular}
\end{table}

The Route~I entries are elementary, which is why they require no distributional assumption.
With $D \ge 0$ and $\tau > 0$, Markov's inequality $\Pr(D > \tau) \le \mathbb{E}[D]/\tau$ follows from $\tau\,\mathbf{1}\{D > \tau\} \le D$.
Applying Markov to $(D - \mathbb{E}[D] + c)^2$ and optimising over $c > 0$ gives the one-sided Chebyshev (Cantelli) bound, for $\tau > \mathbb{E}[D]$,
\begin{equation}\label{eq:cantelli_surrogate}
  \Pr(D > \tau) \;\le\; \frac{\operatorname{Var}[D]}{\operatorname{Var}[D] + (\tau - \mathbb{E}[D])^2}.
\end{equation}
Both are the \emph{tight} bounds attainable from the first one and two moments, respectively~\cite{BertsimasPopescu2005}, and both are attained (for the closed event $D \ge \tau$) by two-point distributions---which is exactly why they ignore, and are loose under, the independent-sum structure.
The exponential (Cram\'er--Chernoff) bound,
\begin{equation}\label{eq:chernoff_surrogate}
  \Pr(D > \tau) \;\le\; \inf_{s > 0} e^{-s\tau} M_D(s)
  \;=\; \inf_{s > 0} \exp\!\Big( -s\tau + \textstyle\sum_t \ln M_{\Delta D_t}(s) \Big),
\end{equation}
follows from Markov applied to $e^{sD}$ and is valid for every fixed $s > 0$~\cite{Chernoff1952}; Bernstein's inequality is the specialisation for independent summands bounded in $[0,b]$, $\Pr(D - \mathbb{E}[D] > a) \le \exp\big( -\tfrac{a^2/2}{\operatorname{Var}[D] + b\,a/3} \big)$ with $a = \tau - \mathbb{E}[D]$~\cite{BoucheronLugosiMassart2013}.
Unlike Markov and Cantelli, these use the additive log-MGF of~\eqref{eq:additive_functionals} and therefore exploit independence.

\paragraph{Per-mission inputs do not flatten the missions.}
Collapsing the summands to a single worst-case support $b = \max_j b_{\ell j}$ would make Bernstein treat every mission as the most damaging one, erasing the differences the schedule exploits; this flattening is an artefact of the closed form, not of Route~I.
Kept per term, each bound retains every mission's own support and tail.
Hoeffding reads $\Pr(D - \mathbb{E}[D] > a) \le \exp\!\big( -2a^2 / \sum_j n_{ij} b_{\ell j}^2 \big)$ with $a = \tau - \mathbb{E}[D]$, where mission~$j$ enters only through its own $b_{\ell j}^2$, and the Chernoff bound~\eqref{eq:chernoff_surrogate} becomes the count-weighted reliability constraint
\begin{equation}\label{eq:chernoff_permission}
  \Pr(D > \tau) \;\le\; \inf_{s > 0} \exp\!\Big( -s\tau + \sum_j n_{ij}\,\psi_{\ell j}(s) \Big),
\end{equation}
in which each mission contributes its own cumulant generating function $\psi_{\ell j}$.
For a fixed~$s$ the exponent is linear in the assignment counts, so~\eqref{eq:chernoff_permission} is a linear reliability constraint; the inner infimum couples~$s$ to the counts---for a fixed schedule it is a one-dimensional convex problem whose minimiser $s^\star$ solves $\sum_j n_{ij}\,\psi_{\ell j}'(s) = \tau$ (the exponentially tilted mean of~$D$ equals~$\tau$), but since $s^\star$ depends on the counts one either fixes a conservative~$s$ (any $s > 0$ remains a valid bound) or sweeps a small grid.

Bernstein's inequality is the one Route~I bound that resists this per-mission treatment, and the reason is the same coupling seen from the other side: its closed form requires a \emph{single} support $b_\ell = \max_j b_{\ell j}$ in the correction term (only the variance stays per-mission).
Explicitly, with $\kappa = \ln(1/\varepsilon)$, the reliability constraint $\Pr(D > \tau) \le \varepsilon$ becomes
\begin{equation}\label{eq:bernstein_constraint}
  \sum_j n_{ij}\,\mu_{\ell j}
  \;+\; \frac{\kappa\, b_\ell}{3}
  \;+\; \sqrt{ \Big( \tfrac{\kappa\, b_\ell}{3} \Big)^2 + 2\kappa \sum_j n_{ij}\, v_{\ell j} }
  \;\le\; \tau,
\end{equation}
a constraint whose left-hand side is \emph{concave} in the assignment counts (equivalently $(\kappa b_\ell/3)^2 + 2\kappa\sum_j n_{ij} v_{\ell j} \le (\tau - \kappa b_\ell/3 - \sum_j n_{ij}\mu_{\ell j})^2$ with the bracket non-negative): by~\eqref{eq:additive_functionals} the mean and the variance are \emph{affine} in the counts, so the square root is concave and the feasible set is reverse-convex rather than second-order-cone-representable---in the classical SOCP chance constraint the variance is a convex quadratic of the decisions, whereas here independence across days makes it affine, which flips the curvature.
For the MILP this is a feature rather than an obstacle: every tangent line to the concave square root \emph{over}-estimates it, so replacing the square root in~\eqref{eq:bernstein_constraint} by a fixed tangent---or by the piecewise-linear pointwise minimum of several tangents, encoded exactly as a PWL equality (SOS2)---yields linear constraints in the counts that imply~\eqref{eq:bernstein_constraint} and therefore remain conservative.
Retaining a per-mission support $b_{\ell j}$ in the correction is possible only by keeping the per-term cumulant sum $\sum_j n_{ij}\,(s^2 v_{\ell j}/2)/(1 - s\,b_{\ell j}/3)$ and minimising over~$s$---which is exactly the Chernoff/Bennett bound~\eqref{eq:chernoff_permission}.
The $s$-free closed form~\eqref{eq:bernstein_constraint} and per-mission support resolution therefore cannot be obtained simultaneously.

\paragraph{Guaranteed bound versus estimate.}
For~\eqref{eq:chernoff_permission} to be a true bound, the per-mission ingredient must be an \emph{upper} bound on $\psi_{\ell j}$.
Two per-mission choices suffice and keep the support local: the mean-and-support bound $\psi_{\ell j}(s) \le \ln\!\big( 1 + (\mu_{\ell j}/b_{\ell j})(e^{s b_{\ell j}} - 1) \big)$, attained by the two-point law on $\{0, b_{\ell j}\}$; or the empirical cumulant generating function $\widehat\psi_{\ell j}(s) = \ln\!\big( \tfrac{1}{m} \sum_{r=1}^{m} e^{s \Delta D^{(j)}_r} \big)$ built from the $m$ recorded executions of mission~$j$, which is an \emph{estimate}, not a bound, and so must carry a one-sided confidence margin.
Repeating a mission tightens both ingredients---more executions sharpen $\widehat\psi_{\ell j}$, $\mu_{\ell j}$ and $b_{\ell j}$ and shrink the margin---which is the precise sense in which undergoing a mission many times refines its tail description.

The choice among the Route~I bounds depends on how deep into the tail the constraint sits.
In the regime relevant to a reliability constraint with small~$\varepsilon$---several standard deviations above $\mathbb{E}[D]$---the exponential bounds are markedly tighter: in a representative numerical experiment (non-negative increments bounded in $[0,1]$, $k \in \{30,100\}$ summands), at $\varepsilon \sim 10^{-4}$ the Chernoff bound stayed within roughly one order of magnitude of the empirical tail, whereas the Cantelli bound over-stated it by more than two orders; Bernstein fell between them and overtook Cantelli only as the tail deepened and the number of summands grew.\footnote{Recorded as a numerical illustration, not a proof; the test code and output are kept in the verification workspace. Asymptotically the gap is structural: Cantelli decays only polynomially, as $\operatorname{Var}[D]/(\tau-\mathbb{E}[D])^2$, whereas~\eqref{eq:chernoff_surrogate} decays exponentially in the deviation.}
Cantelli, though the simplest, is therefore best reserved for short accumulation windows or as a coarse always-valid fallback.

Route~II yields not a bound but an explicit surrogate law.
Matching only the first two moments of a convolution-closed family (Gamma, with shape $\alpha = \mathbb{E}[\Delta D]^2/\operatorname{Var}[\Delta D]$ and rate $\beta = \mathbb{E}[\Delta D]/\operatorname{Var}[\Delta D]$; or inverse Gaussian) is convenient---the sum stays in the family under the closure conditions of Table~\ref{tab:discrete_models}---but is \emph{not} conservative: a moment-matched law generally crosses the empirical survival function rather than dominating it.
Conservatism is recovered by using more of the empirical information than its moments.
Fix a common rate $\beta^\star$ and, per mission, inflate the shape $\alpha_j$ until the surrogate survival dominates the empirical one, $\overline{\mathrm{Ga}}(x;\alpha_j,\beta^\star) \ge \widehat{\bar F}_j(x)$ for all $x$---always achievable, since $\overline{\mathrm{Ga}}(x;\alpha,\beta^\star) \to 1$ as $\alpha \to \infty$.
By Proposition~\ref{prop:st_closure} the accumulated surrogate is then $\widetilde{D} = \mathrm{Ga}(A,\beta^\star)$ with cumulative shape $A = \sum_{t,j} x_{ijt}\,\alpha_j$ linear in the schedule, the reliability constraint reduces to the scalar condition $A \le A^\star(\varepsilon)$, and the loop constraint to $A_{2H} \le A_H$ exactly as in Equation~\eqref{eq:loop_gamma}.
The price is the slack introduced by the shared rate and the dominance margin.

In summary, for the one-sided reliability constraint we use the count-weighted Chernoff bound~\eqref{eq:chernoff_permission}---which keeps each mission's own tail---or its per-term Hoeffding specialisation when only per-mission means and supports are trusted; when an explicit surrogate law is required---notably for the loop constraint, a comparison of two survival functions---we use the dominating shared-rate Gamma of the last row, which lands on the tractable Gamma machinery of Section~\ref{sec:stoch_orders}.
The irreducible trade-off is that the bounds assume nothing about the shape but yield only a tail bound, whereas the dominating Gamma yields a propagatable law at the cost of a family commitment and some conservatism.

\subsubsection{The loop constraint under the Route-I bounds}
\label{sec:loop_routeI}

The loop constraint~\eqref{eq:loop} requires $\Pr(D_{2H} > \tau) \le \Pr(D_H > \tau)$ so that concatenating the schedule keeps the failure probability below~$\varepsilon$ indefinitely.
Under Route~I we never observe these probabilities; we have only an upper bound $U_k \ge \Pr(D_k > \tau)$ from the count-weighted Chernoff bound~\eqref{eq:chernoff_permission} or the Bernstein constraint~\eqref{eq:bernstein_constraint}.
The tempting translation---replace both sides by their bounds and impose $U_{2H} \le U_H$---is \emph{not} sufficient for repeatability, and it is worth seeing precisely why.

\paragraph{Why $U_{2H} \le U_H$ does not chain.}
The bound $U$ is a many-to-one function of the state's parameters, and its level sets are not invariant under the dynamics.
Bernstein's bound, for instance, depends on $\big(\mathbb{E}[D],\operatorname{Var}[D]\big)$ through a single scalar, so $U_{2H} = U_H$ can hold with $\big(\mathbb{E}[D_{2H}],\operatorname{Var}[D_{2H}]\big) \ne \big(\mathbb{E}[D_H],\operatorname{Var}[D_H]\big)$---a larger mean offset by a smaller variance, say.
One further period of the dynamics maps these two distinct parameter points to \emph{different} bound values, so $U_{3H} > U_{2H} = U_H$ is possible: the certified bound can creep upward over repetitions, and the single-period inequality does not induct.\footnote{A concrete instance: $\tau = 10$, $b_\ell = 3$, $\rho = 0.5$, repeated block $=$ one repair day followed by mission days with $\sum_j n_j\mu_{\ell j} = 3$ and $\sum_j n_j v_{\ell j} = 3.99$; from $\bigl(\mathbb{E}[D_H],\operatorname{Var}[D_H]\bigr) = (4,9)$ the Cantelli bound gives $U_H = 0.2000$ and $U_{2H} = 0.1997 \le U_H$, yet $U_{3H} = 0.2151 > U_H$, with $U_{mH} \to 0.2495$. An analogous Bernstein instance (block variance $3.15$) violates at $m=3$ as well. The moment condition~\eqref{eq:loop_moments} rejects both, since $\mathbb{E}[D_{2H}] = 5 > 4$. Test scripts are kept in the verification workspace.}
Adding a fixed margin, $U_{2H} \le U_H - c$, does not repair this, because the defect is the non-invariance of the level sets, not a constant gap.

\paragraph{The certificate stays valid at every period.}
This is not an undetectable worsening.
The functionals Route~I actually tracks---the moments $\mathbb{E}[D_k],\operatorname{Var}[D_k]$, or the cumulant generating function---are exact functionals of the true state and are \emph{closed} under the dynamics~\eqref{eq:dyn_full}: increments add them, imperfect repair scales them, replacement resets them.
Consequently $U_k$ is a valid bound at \emph{every} period, $\Pr(D_{kH} > \tau) \le U_{kH}$, and the true probability can exceed~$\varepsilon$ only if $U_{kH}$ itself does.
Repeatability therefore needs a loop condition that forces $U_{kH} \le \varepsilon$ for all~$k$, which means imposing it on a statistic that both determines $U$ and propagates monotonically---not on the lossy scalar~$U$.

\paragraph{Looping on the moments (the native condition).}
Impose the loop on the exactly tracked moments,
\begin{equation}\label{eq:loop_moments}
  \mathbb{E}[D_{2H}] \;\le\; \mathbb{E}[D_H]
  \qquad\text{and}\qquad
  \operatorname{Var}[D_{2H}] \;\le\; \operatorname{Var}[D_H].
\end{equation}
Under the dynamics, the mean and variance evolve by an increment ($\mathbb{E} \mathrel{+}= \sum_j n_{ij}\mu_{\ell j}$, $\operatorname{Var} \mathrel{+}= \sum_j n_{ij} v_{\ell j}$), by imperfect repair ($\mathbb{E} \mapsto (1-\rho)\,\mathbb{E}$, $\operatorname{Var} \mapsto (1-\rho)^2\operatorname{Var}$), or by replacement (reset to the constants of $D^{\mathrm{new}}$); each operation is non-decreasing in $\mathbb{E}$ and in $\operatorname{Var}$ separately, so the two moments evolve as independent monotone recursions.
Componentwise domination therefore chains: writing $F$ for the once-per-period (second-half) map, \eqref{eq:loop_moments} states $F(\,\cdot\,)$ does not increase either moment, whence $\big(\mathbb{E}[D_{(m+1)H}],\operatorname{Var}[D_{(m+1)H}]\big) \le \big(\mathbb{E}[D_{mH}],\operatorname{Var}[D_{mH}]\big)$ for every~$m \ge 1$.
Because the Bernstein bound~\eqref{eq:bernstein_constraint}---and likewise the Cantelli bound~\eqref{eq:cantelli_surrogate}---is increasing in both $\mathbb{E}[D]$ and $\operatorname{Var}[D]$, the certified risk is then non-increasing across periods, $U_{kH} \le U_H \le \varepsilon$, and the schedule is repeatable.
Both inequalities in~\eqref{eq:loop_moments} are linear in the tracked moment states $\bigl(\mathbb{E}[D_k],\operatorname{Var}[D_k]\bigr)$, which the model carries as auxiliary variables (the moment recursion couples the maintenance binaries to these states through products, linearised exactly by standard big-$M$ constructions), so the loop adds two linear constraints per component; they are strictly stronger than the single relation $U_{2H} \le U_H$, and that strengthening is exactly what makes the induction valid.

\paragraph{Looping on the cumulant, and why imperfect repair forbids it.}
If reliability is instead certified by a fixed-$s$ Chernoff bound, the loop may be placed on the cumulant,
\begin{equation}\label{eq:loop_cumulant}
  \sum_j n^{(2H)}_{ij}\,\psi_{\ell j}(s_0) \;\le\; \sum_j n^{(H)}_{ij}\,\psi_{\ell j}(s_0),
\end{equation}
where $n^{(k)}_{ij}$ is the post-reset cumulative count of mission~$j$ up to day~$k$.
This chains \emph{provided the repeated block contains only additive increments and full replacement}: there $M_{D}(s_0)$ propagates multiplicatively, $M_{D + \Delta D}(s_0) = M_D(s_0)\,M_{\Delta D}(s_0)$, so it is a self-propagating scalar and~\eqref{eq:loop_cumulant} inducts.
Imperfect (ARD) repair breaks this. With $D^+ = (1-\rho)D^-$,
\begin{equation}\label{eq:mgf_rescale}
  M_{D^+}(s_0) \;=\; \mathbb{E}\big[e^{s_0(1-\rho)D^-}\big] \;=\; M_{D^-}\big((1-\rho)s_0\big):
\end{equation}
the repaired state's transform at $s_0$ is the old transform at the \emph{shifted} argument $(1-\rho)s_0$, which is not determined by $M_{D^-}(s_0)$---the repair slides the state along the transform curve.
The single scalar $M(s_0)$ is then no longer self-contained, the count-weighted cumulant ceases to be additive across the repair, and~\eqref{eq:loop_cumulant} no longer chains.
The moments escape this because scaling acts on each of them by a fixed power, $\mathbb{E}\mapsto(1-\rho)\mathbb{E}$ and $\operatorname{Var}\mapsto(1-\rho)^2\operatorname{Var}$, keeping their recursion closed---which is why~\eqref{eq:loop_moments} is the form to use whenever imperfect repair appears inside the repeated block.

\paragraph{Summary.}
Under Route~I the loop constraint must be imposed on a sufficient, self-propagating statistic, never on the scalar bound~$U$.
The componentwise moment condition~\eqref{eq:loop_moments} is the native choice: two linear constraints, valid for the Bernstein and Cantelli certificates and for every maintenance operator in~\eqref{eq:dyn_full}, imperfect repair included.
The cumulant condition~\eqref{eq:loop_cumulant} is available when the certificate is the fixed-$s$ Chernoff bound and the repeated block is free of imperfect repair.
In either case the guarantee is non-worsening of the \emph{certified} risk, hence $\Pr(D_{kH} > \tau) \le \varepsilon$ for all concatenations---the operative meaning of repeatability for a bound-based model.

\subsubsection{Two refinements of the Route-I constraints: schedule-dependent support and piecewise-linear tightening}
\label{sec:routeI_refinements}

Two identifiable sources of slack remain in the Route-I constraints, and both can be removed with standard mixed-integer constructions at a quantifiable price in binaries and rows.
Neither touches the loop argument of Section~\ref{sec:loop_routeI}: the moment condition~\eqref{eq:loop_moments} acts on the exactly tracked moments and is agnostic to how tightly the certificate is enforced on them.\footnote{Every monotonicity and equivalence claim in this subsection, both worked examples, and the union-versus-intersection geometry are checked numerically; the test scripts and outputs are kept in the verification workspace (\texttt{refine\_schedule\_support}, \texttt{refine\_pwl\_tightening}).}

\paragraph{Schedule-dependent Bernstein support.}
The closed form~\eqref{eq:bernstein_constraint} carries the single constant $b_\ell = \max_j b_{\ell j}$ (joined, once resets enter the accumulation window, by the supports $b^{\mathrm{start}}_\ell$, $b^{\mathrm{new}}_\ell$ of the initial and replacement draws), so a component pays the most damaging mission's support even if the schedule never assigns it that mission.
The effect is not marginal, because the correction $\kappa b_\ell/3$ enters twice at zero variance: even the empty schedule satisfies~\eqref{eq:bernstein_constraint} only if $2\kappa b_\ell/3 \le \tau$.
Concretely, take $\varepsilon = 0.01$ ($\kappa = \ln 100 \approx 4.605$), $\tau = 1$, two missions with supports $b_{\ell 1} = 0.02$ and $b_{\ell 2} = 0.30$, reset supports $0.02$, and mission-1 per-day moments $\mu_{\ell 1} = 0.01$, $v_{\ell 1} = 4\cdot 10^{-6}$.
For a component that never serves mission~2, the constant-$b$ constraint ($\kappa b_\ell/3 = 0.4605$) leaves a mean budget of $0.079$ and admits at most $n_{i1} = 7$ mission-1 days, whereas the same constraint evaluated with the supports actually present ($\kappa b_\ell/3 = 0.0307$) leaves $0.939$ and admits $n_{i1} = 90$---a $12.9\times$ gap attributable entirely to the unused $b_{\ell 2}$; and raising $b_{\ell 2}$ to $0.35$ makes $2\kappa b_\ell/3 = 1.075 > \tau$, so the constant-$b$ model becomes infeasible outright for a component that never touches mission~2.
The tightened constant is $b_\ell(\zeta) = \max\bigl\{b^{\mathrm{start}}_\ell,\, b^{\mathrm{new}}_\ell,\, \max_{j:\,\zeta_{ij}=1} b_{\ell j}\bigr\}$, where $\zeta_{ij} \in \{0,1\}$ flags the missions member~$i$ actually serves ($\zeta_{ij} \ge x_{ijt}$ for all~$t$; only this forcing direction is needed): between resets the state unrolls as a $[0,1]$-weighted sum of one reset draw and of increments of \emph{served} missions only, so $b_\ell(\zeta)$ bounds every summand and Bernstein applies verbatim, at every step and under every concatenation of the schedule.
What makes $b_\ell(\zeta)$ implementable without nonlinear max-tracking is a monotonicity observation: the certified value $\mu + c + \sqrt{c^2 + 2\kappa v}$ has $c$-derivative $1 + c/\sqrt{c^2 + 2\kappa v} \in [1,2]$, so it is increasing in the support constant, and the Bernstein row written for a larger constant implies every row written for a smaller one.
Consequently, imposing the row with the per-mission constant $c_j = \kappa\max\{b_{\ell j}, b^{\mathrm{start}}_\ell, b^{\mathrm{new}}_\ell\}/3$ \emph{conditionally on} $\zeta_{ij} = 1$, together with one unconditional row at $c_0 = \kappa\max\{b^{\mathrm{start}}_\ell, b^{\mathrm{new}}_\ell\}/3$, is exactly equivalent to the single row with the schedule-dependent maximum: the active row with the largest constant dominates the rest, and that largest constant is precisely $\kappa b_\ell(\zeta)/3$.
The price is one binary per member--mission pair ($F \cdot M$ in total), their linking rows, and per-mission conditional certificate rows, all linear in the new material: the tangent-surrogate rows attach to $\zeta_{ij}$ by indicator constraints or by big-$M$ (a uniform valid constant is $2\tau(\tau - c_0)$), whereas the exact quadratic rows must use big-$M$, because indicator constraints admit only linear bodies.
Two side effects deserve record.
First, big-$M$ constants derived from the largest admissible variance must be recomputed from $c_0$: the feasible region has grown, and a bound inherited from the old constant $b_\ell$ would cut it.
Second, the feasibility precondition $2\kappa b_\ell/3 < \tau$ localises: it is required only of $c_0$, while a mission with $2c_j \ge \tau$ simply forces $\zeta_{ij} = 0$---that mission is automatically never assigned to that component, instead of rendering the whole model inconsistent.
Deliberately \emph{not} captured: the repair weights further shrink each summand's support to $w_t\,b_{\ell j}$, but exploiting this requires a running maximum over schedule-dependent products of repair factors and mission indicators---a genuinely nonlinear refinement; the window-level $\zeta$ captures the first-order gap, namely missions never served at all.

\paragraph{Piecewise-linear tightening of the tangent surrogates.}
The linear surrogate obtained by fixing one tangent (the device noted after~\eqref{eq:bernstein_constraint}) is conservative in a way that grows quadratically with the mean.
Cantelli's bound~\eqref{eq:cantelli_surrogate} is at most $\varepsilon$ iff $\mathbb{E}[D] + \kappa_C\sqrt{\operatorname{Var}[D]} \le \tau$ with $\kappa_C = \sqrt{(1-\varepsilon)/\varepsilon}$; writing $(\mu, v)$ for the tracked moments, the exact feasible set is $\kappa_C^2\, v \le f(\mu) := (\tau - \mu)^2$, and the tangent at $\mu_0 = 0$ gives the linear row $\kappa_C^2\, v \le \tau^2 - 2\tau\mu$.
Since $f(\mu) - (\tau^2 - 2\tau\mu) = \mu^2$, the surrogate loses exactly $\mu^2/\kappa_C^2$ of admissible variance, reaches zero at $\mu = \tau/2$, and excludes \emph{every} point with $\mu > \tau/2$---a deterministic component included---while the exact set still admits $v = (\tau - \mu)^2/\kappa_C^2 > 0$ up to $\mu = \tau$.
The refinement takes tangents at several points $\mu_r$: each row $\kappa_C^2\, v \le T_r(\mu) := (\tau - \mu_r)(\tau + \mu_r - 2\mu)$ is individually conservative, because tangents lie below the convex parabola.
The directionality is the critical point: imposing all rows \emph{simultaneously} realises their intersection $\kappa_C^2 v \le \min_r T_r(\mu)$, which is \emph{more} conservative than any single tangent; the tightening is the \emph{union} of the single-tangent regions, $\kappa_C^2 v \le \max_r T_r(\mu)$---a disjunction.
Two encodings realise it: (i)~one selector binary per tangent, $\sum_r \lambda_r = 1$, with indicator (or big-$M$) rows $\lambda_r = 1 \Rightarrow \kappa_C^2 v \le T_r(\mu)$; or (ii)~an auxiliary $y$ tied by a piecewise-linear \emph{equality} (SOS2 $\lambda$-formulation) to the pointwise minimum of the tangents to $\sqrt{v}$ at $v_r = (\tau - \mu_r)^2/\kappa_C^2$, followed by the single linear row $\mu + \kappa_C\, y \le \tau$.
The two describe the same set: the $\mu$-side tangent at $\mu_r$ and the $v$-side tangent at $v_r$ bound the identical half-plane $\kappa_C^2 v \le (\tau - \mu_r)(\tau + \mu_r - 2\mu)$.
The tempting shortcuts both fail: requiring $y \ge$ every tangent realises their pointwise \emph{maximum} (a region strictly inside the single-tangent one, worse than not refining), and requiring $y \le$ every tangent lets $y$ under-run $\sqrt{v}$ and destroys the certificate; likewise the equality must interpolate the vertices of the tangent envelope---consecutive tangents meet at $v = \sqrt{v_r v_{r+1}}$ with value $(\sqrt{v_r} + \sqrt{v_{r+1}})/2$, above the curve---not points on the curve, whose chords under-estimate the concave square root.
The recovery is rapid, and exactly quantifiable because $f$ is quadratic: tangent~$r$'s gap is $f(\mu) - T_r(\mu) = (\mu - \mu_r)^2$, so the union's residual gap is $\min_r (\mu - \mu_r)^2$---quadratic in the distance to the nearest tangency, with worst case $(\Delta/2)^2$ for spacing $\Delta$, so doubling the tangents quarters it.
With $\tau = 1$, $\varepsilon = 0.01$ ($\kappa_C^2 = 99$) and tangents at $\mu_r \in \{0, 0.25, 0.5, 0.75\}$, the admissible variance recovers as follows ($v \times 10^{3}$):
\begin{center}
\footnotesize
\begin{tabular}{@{}lccccc@{}}
\toprule
$\mu$ & $0$ & $0.2$ & $0.4$ & $0.5$ & $0.6$ \\
\midrule
exact $(\tau-\mu)^2/\kappa_C^2$ & $10.10$ & $6.46$ & $3.64$ & $2.53$ & $1.62$ \\
single tangent ($\mu_0 = 0$)    & $10.10$ & $6.06$ & $2.02$ & $0$    & infeasible \\
four tangents (union)           & $10.10$ & $6.44$ & $3.54$ & $2.53$ & $1.52$ \\
\bottomrule
\end{tabular}
\end{center}
and the union recovers $94\%$ of the feasible-region area the single tangent forfeits.
The construction transfers verbatim to Bernstein---the parabola $(\tau - \kappa b_\ell/3 - \mu)^2$ obeys the same gap law, and the tangents to the concave $\sqrt{(\kappa b_\ell/3)^2 + 2\kappa v}$ over-estimate it globally---and composes soundly with the schedule-dependent support above: each conditional row carries its own mission constant and remains conservative for it.

\subsubsection{Random fatigue life in the Palmgren--Miner rule}
\label{sec:miner_randomN}

The rainflow pipeline of Section~\ref{sec:cons_surrogates} measures the law of the per-mission increment $\Delta D^{(j)}$ directly and nonparametrically.
A complementary pipeline keeps the rainflow \emph{counts} deterministic and makes the fatigue life $N_f$ in Miner's rule~\eqref{eq:miner} the random object: one fits a scatter family of Table~\ref{tab:sn_scatter} to constant-amplitude test data or in-service failure records, tying the location to the stress amplitude through Basquin's law~\eqref{eq:basquin}---for the log-normal, $\ln N_f \sim \mathcal{N}(\mu(\sigma_a),\sigma^2)$ with $\mu(\sigma_a) = \ln C - m \ln \sigma_a$ and a stress-independent $\sigma$; for the Weibull, scale $\lambda(\sigma_a) = C\sigma_a^{-m}$ and common shape $k$; run-outs enter the likelihood as censored observations---and then \emph{derives} the increment from the mission's rainflow histogram $\{n^{(j)}_c\}$ as $\Delta D^{(j)} = \sum_c n^{(j)}_c / N_f(\sigma_c)$.
The retrieval is the same in spirit---collect data, approximate a distribution---but dual in object and commitment: the empirical route is distribution-free and captures load and material variability jointly, whereas the random-$N_f$ route is parametric, pools information across stress levels, extrapolates to unobserved amplitudes through~\eqref{eq:basquin}, and separates the load spectrum (the counts) from the material scatter (the law of $N_f$).

\paragraph{Inverse moments of the scatter families.}
Everything the framework consumes downstream is a moment of $1/N_f$.
For the three families of Table~\ref{tab:sn_scatter} and $r > 0$,
\begin{equation}\label{eq:miner_invmom}
  \mathbb{E}\big[N_f^{-r}\big] \;=\;
  \begin{cases}
    \exp\big({-r\mu + r^2\sigma^2/2}\big), & \ln N_f \sim \mathcal{N}(\mu,\sigma^2) \quad \text{(finite for all $r$)},\\[3pt]
    \lambda^{-r}\,\Gamma(1 - r/k), & N_f \sim \mathrm{Wei}(k,\lambda) \quad \text{(finite iff $r < k$)},\\[3pt]
    e^{-r\mu}\,\Gamma(1 - r\beta), & \ln N_f \sim \mathrm{Gumbel}_{\min}(\mu,\beta) \quad \text{(finite iff $r\beta < 1$)}.
  \end{cases}
\end{equation}
The first line is the normal moment generating function at argument $-r$; the second follows from the substitution $u = (x/\lambda)^k$, which turns $\mathbb{E}[N_f^{-r}]$ into $\lambda^{-r}\int_0^\infty u^{-r/k}e^{-u}\,\mathrm{d}u$, divergent \emph{at the origin} when $r \ge k$; the third writes $\ln N_f = \mu + \beta G$ with $\Pr(G \le g) = 1 - e^{-e^{g}}$, whose transform $\mathbb{E}[e^{tG}] = \int_0^\infty u^{t}e^{-u}\,\mathrm{d}u = \Gamma(1+t)$ (substitute $u = e^g$) is finite iff $t > -1$.
In particular $\operatorname{Var}[1/N_f]$ equals $e^{-2\mu+\sigma^2}(e^{\sigma^2}-1)$, $\lambda^{-2}\big[\Gamma(1-2/k)-\Gamma(1-1/k)^2\big]$, and $e^{-2\mu}\big[\Gamma(1-2\beta)-\Gamma(1-\beta)^2\big]$ respectively, matching the increment moments already recorded for the log-normal below Table~\ref{tab:sn_scatter}.

\begin{remark}[Two tail behaviours, one real restriction]\label{rem:miner_tails}
The Gumbel-minimum log-life is \emph{exactly} a Weibull life: $\Pr(N_f > x) = \exp\big(-e^{(\ln x - \mu)/\beta}\big) = \exp\big(-(x/e^{\mu})^{1/\beta}\big)$, i.e.\ $\mathrm{Wei}(k = 1/\beta,\ \lambda = e^{\mu})$, and the existence conditions in~\eqref{eq:miner_invmom} map onto each other.
The three families of Table~\ref{tab:sn_scatter} therefore realise two tail behaviours for $1/N_f$: log-normal (all polynomial moments finite) and Fr\'echet-tailed with index $k$ (moments capped at order $k$).
The latter is a genuine restriction, not a formality: the mean of the increment requires $k > 1$ and the variance---the input of the Cantelli certificate~\eqref{eq:cantelli_surrogate}---requires $k > 2$, so the fitted shape must be checked before the machinery is invoked; for $k \in (1,2]$ only the Markov bound survives, and for $k \le 1$ the untruncated model has an infinite mean increment and cannot enter the framework at all.
Moreover, for \emph{all three} untruncated families $1/N_f$ has unbounded support and an infinite moment generating function, so the exponential certificates---Chernoff~\eqref{eq:chernoff_permission}, Bernstein~\eqref{eq:bernstein_constraint}, Hoeffding---are inadmissible; they are restored only by a physically defended lower life bound $N_{\min} > 0$ (truncated family), which yields the per-mission support $b_{\ell j} = n_j / N_{\min}$ and requires the moments to be recomputed under the truncated law (Gaussian-tail and incomplete-gamma closed forms).
\end{remark}

\paragraph{What ``random $N_f$'' means: two interpretations.}
The marginal law of $N_f$ does not determine the model; the joint law of the scatter \emph{across days} must be chosen, and the two extremes lead to different mathematics.

\emph{(I) Block-independent scatter.}
Each execution of mission~$j$ draws a fresh, independent life---appropriate when cycle-to-cycle and day-to-day variability dominates and the component does not carry one fixed $N_f$.
With a single dominant stress level per mission (for a binned spectrum the means add over bins, while the variance depends on the within-mission dependence assumption and is computed from~\eqref{eq:miner_invmom} in the same way), the per-mission inputs are
\begin{equation}\label{eq:miner_routeI}
  \mu_{\ell j} \;=\; n_j\,\mathbb{E}\big[N_f^{-1}\big],
  \qquad
  v_{\ell j} \;=\; n_j^2\,\operatorname{Var}\big[N_f^{-1}\big],
\end{equation}
and independence across days is precisely the hypothesis of Equation~\eqref{eq:additive_functionals}: the Cantelli certificate~\eqref{eq:cantelli_surrogate} and the moment loop~\eqref{eq:loop_moments} apply \emph{verbatim}, with parametric instead of empirical per-mission moments.
This is the interpretation the rainflow route of Section~\ref{sec:cons_surrogates} already embodies; the random-$N_f$ model is then nothing but an alternative input pipeline for the same formulation, subject to the admissibility conditions of Remark~\ref{rem:miner_tails}.

\emph{(II) Unit-specific scatter.}
The component carries \emph{one} realisation of the scatter for its entire life.
Under the stress-independent scale/shape of the fitted families this is a single multiplicative factor,
\begin{equation}\label{eq:miner_unit}
  N_{f,j} \;=\; A\,\bar N_j, \qquad A > 0 \ \text{drawn once at installation},
\end{equation}
with $\bar N_j$ the median (log-normal, Gumbel: $e^{\mu_j}$) or characteristic (Weibull: $\lambda_j$) life and $A \sim F_A$ the unit factor ($\ln A \sim \mathcal{N}(0,\sigma^2)$, $A \sim \mathrm{Wei}(k,1)$, or $A = e^{\beta G}$ respectively).
This is exactly the ``random critical damage sum'' of source~(ii) above: dividing the Miner sum by $A$ is the same as comparing a deterministic sum to a random capacity.
The damage state becomes $D_k = M_k / A$, where $M_k$ is the \emph{deterministic} Miner state with dynamics mirroring~\eqref{eq:dyn_full}: a mission day adds $\sum_j x_{ijk}\, n_j/\bar N_j$, ARD repair maps $M \mapsto (1-\rho)M$, and replacement resets $M$ to zero \emph{and redraws} $A$ (a new physical component carries a fresh factor with the same law).
For continuous $F_A$ the reliability constraint is then \emph{exact}, deterministic, and linear in the assignment counts:
\begin{equation}\label{eq:miner_quantile}
  \Pr(D_k \ge 1) \;=\; \Pr(A \le M_k) \;=\; F_A(M_k) \;\le\; \varepsilon
  \quad\Longleftrightarrow\quad
  M_k \;\le\; q_\varepsilon \;:=\; F_A^{-1}(\varepsilon),
\end{equation}
with closed-form quantiles $q_\varepsilon = e^{\sigma\Phi^{-1}(\varepsilon)}$ (log-normal), $q_\varepsilon = (-\ln(1-\varepsilon))^{1/k}$ (Weibull), $q_\varepsilon = (-\ln(1-\varepsilon))^{\beta}$ (Gumbel-minimum); indeed no family is needed at all beyond one scalar---an empirical $\varepsilon$-quantile of $A$ with a confidence margin serves equally.
Rewriting $M_k \le q_\varepsilon$ as $\sum_j n_{ij}/(q_\varepsilon \bar N_j) \le 1$ recovers the classical design practice of running deterministic Miner on the $\varepsilon$-percentile lives $N^{(\varepsilon)}_j = q_\varepsilon \bar N_j$ (these are exactly the $\varepsilon$-quantiles of $N_{f,j}$)---here as an \emph{exact} chance constraint rather than a heuristic, in the per-step sense of~\eqref{eq:abs_rel}.
The loop constraint collapses likewise: $F_A$ is strictly increasing on $(0,\infty)$, so~\eqref{eq:loop} is equivalent to the single linear condition
\begin{equation}\label{eq:miner_loop_det}
  M_{2H} \;\le\; M_H,
\end{equation}
and---unlike the lossy scalar bound of Section~\ref{sec:loop_routeI}---it \emph{chains}: $M$ is a sufficient statistic (the failure probability is the fixed increasing function $F_A$ of the scalar state), and the second-half map $F$ is a composition of $m \mapsto m + c$ ($c \ge 0$), $m \mapsto (1-\rho)m$ and $m \mapsto 0$, each non-decreasing, so from $M_{2H} = F(M_H) \le M_H$ induction gives $M_{(m+1)H} = F(M_{mH}) \le F(M_{(m-1)H}) = M_{mH}$ for all $m$; the same monotonicity applied to the partial flows dominates every within-period trajectory by the period-two trajectory, so imposing~\eqref{eq:miner_quantile} on all $k = 1,\dots,2H$ extends the per-step reliability to every concatenation, replacements (which redraw $A$) included.

\begin{remark}[The dichotomy is not a conservatism knob]\label{rem:miner_dichotomy}
Interpretations~(I) and~(II) can share the \emph{same} per-day marginal law and differ only in dependence---zero versus perfect correlation of the scatter across days---yet they answer differently.
Means coincide, while by Cauchy--Schwarz $\operatorname{Var}_{\mathrm{II}}/\operatorname{Var}_{\mathrm{I}} = \big(\sum_t m_t\big)^2 / \sum_t m_t^2$ (the number of days, for equal daily increments $m_t$): under~(II) the dispersion never averages out.
In the small-$\varepsilon$ regime~(II) is by far the heavier-tailed---in a representative experiment (log-normal factor, $\sigma = 0.5$, fifty days, median Miner sum $0.6$) the failure probability was $0.15$ under~(II) and $\sim 10^{-6}$ under~(I)---while at high accumulated damage the ordering flips, so neither interpretation is uniformly conservative.\footnote{Test code and output (\texttt{miner\_a}--\texttt{miner\_c}) are kept in the verification workspace, covering the closed forms~\eqref{eq:miner_invmom}, the divergent sample variance at Weibull $k \le 2$, the exactness of~\eqref{eq:miner_quantile} along a schedule with repair and replacement, the chaining of~\eqref{eq:miner_loop_det}, and the size of the (I)/(II) gap.}
The choice is an identification question: scatter traced to specimen-to-specimen variability (material batch, manufacturing) is unit-specific~(II); scatter from load-history randomness within a mission is block-independent~(I); rainflow-counted in-service data from a fleet mixes both, and a between-unit/within-unit variance decomposition of the recorded increments estimates the split.
When both sources matter one writes $N_f = A\,\widetilde N_f$ with a unit factor and block-independent residuals; conservative composites exist---e.g.\ $\varepsilon$-splitting, $\Pr(\text{fail}) \le \delta + \Pr(\text{fail} \mid A = a_\delta)$ with $a_\delta$ the $\delta$-quantile of $A$, after which the increments are again independent across days with moments scaled by $a_\delta^{-1}$ and $a_\delta^{-2}$ and the Route-I machinery runs on the residual budget $\varepsilon - \delta$---but their calibration is left as an outlook.
\end{remark}

\paragraph{Implications for the abstract formulation.}
Interpretation~(I) requires no change to Problem~\eqref{eq:abstract}: it supplies the per-mission moments~\eqref{eq:miner_routeI} to the same Route-I pipeline that consumes the rainflow-measured ones, and inherits its certificates and loop condition unchanged.
Interpretation~(II) defines a strictly simpler model family: the stochastic state collapses to a deterministic scalar $M_{i\ell k}$ with linear dynamics, the chance constraint becomes the exact linear inequality~\eqref{eq:miner_quantile}, and the loop becomes the linear condition~\eqref{eq:miner_loop_det}---no moment tracking, no tail-bound slack, and no big-$M$ linearisation of a variance recursion---with the entire stochastic description compressed into one quantile $q_\varepsilon$ per component.
It is therefore a natural candidate for a dedicated (deterministic-equivalent) model type alongside the process-based ones, with the caveat of Remark~\ref{rem:miner_dichotomy}: which interpretation matches a given component is an empirical question, not a formulation preference.

\subsection{Stochastic orderings}
\label{sec:stoch_orders}

Comparing the degradation state of a component at different points in the schedule requires a notion of ``one random variable being larger than another.''
Several stochastic orders formalise this idea with increasing strength~\cite{ShakedShanthikumar2007,WikiStochDom}.
Let $X$ and $Y$ be two non-negative random variables with CDFs $F_X$, $F_Y$, survival functions $\bar F_X = 1-F_X$, $\bar F_Y = 1-F_Y$, and densities $f_X$, $f_Y$ (when they exist).

\begin{definition}[Usual stochastic order]\label{def:st}
  $X \le_{\mathrm{st}} Y$ if $F_X(t) \ge F_Y(t)$ for all $t \in \mathbb{R}$,
  or equivalently $\Pr(X > t) \le \Pr(Y > t)$ for all~$t$.
\end{definition}

\begin{definition}[Hazard rate order]\label{def:hr}
  $X \le_{\mathrm{hr}} Y$ if $\bar F_Y(t)/\bar F_X(t)$ is non-decreasing in $t$ over $\{t : \bar F_X(t) > 0\}$.
  Equivalently, the hazard rates satisfy $r_X(t) \ge r_Y(t)$ for all~$t$, where $r(t) = f(t)/\bar F(t)$.
\end{definition}

\begin{definition}[Likelihood ratio order]\label{def:lr}
  $X \le_{\mathrm{lr}} Y$ if $f_Y(t)/f_X(t)$ is non-decreasing in $t$ over the union of their supports.
\end{definition}

These three orders are progressively stronger~\cite[Ch.~1]{ShakedShanthikumar2007}.
The implications are strict---none of the converses holds in general:
\begin{equation}\label{eq:order_implications}
  X \le_{\mathrm{lr}} Y
  \;\Longrightarrow\;
  X \le_{\mathrm{hr}} Y
  \;\Longrightarrow\;
  X \le_{\mathrm{st}} Y
  \;\Longrightarrow\;
  \mathbb{E}[X] \le \mathbb{E}[Y].
\end{equation}

\begin{remark}[Connection to the loop constraint]
The loop constraint~\eqref{eq:loop} requires $\Pr(D_{i\ell,2H} > \tau) \le \Pr(D_{i\ell,H} > \tau)$, which is a \emph{pointwise} comparison of the survival functions at the single threshold~$\tau$.
This is strictly weaker than $D_{i\ell,2H} \le_{\mathrm{st}} D_{i\ell,H}$, which demands the inequality at every point.
Establishing the usual stochastic order is therefore a convenient sufficient condition for the loop constraint, though not a necessary one.
\end{remark}

\subsubsection{Stochastic ordering within each degradation family}

For the fleet management problem, the relevant comparison is between two damage states of the same component---for instance the state at mid-horizon versus end-of-horizon---both belonging to the same parametric family.
Tables~\ref{tab:ordering_lr}--\ref{tab:ordering_st} give, for each family from Table~\ref{tab:deg_models}, the necessary and sufficient conditions for each of the three orderings.
An entry marked ``---'' indicates that no simple closed-form characterisation is known to the author; for the inverse Gaussian rows the $\le_{\mathrm{lr}}$ conditions remain a valid sufficient criterion via~\eqref{eq:order_implications}, while for the compound Poisson rows see the discussion below.

\begin{table}[htbp]
\centering
\caption{Conditions for $X \le_{\mathrm{lr}} Y$ (likelihood ratio order) within each family.  Conditions are necessary and sufficient except where noted.}
\label{tab:ordering_lr}
\small
\begin{tabular}{@{}p{2.1cm}p{3.8cm}p{4.0cm}p{4.0cm}@{}}
\toprule
\textbf{Family} & $X$ & $Y$ & \textbf{Condition} \\
\midrule
Gamma
  & $\mathrm{Ga}(\alpha_1,\,\beta_1)$
  & $\mathrm{Ga}(\alpha_2,\,\beta_2)$
  & $\alpha_1 \le \alpha_2$ and $\beta_1 \ge \beta_2$ \\[4pt]
Gaussian
  & $\mathcal{N}(\mu_1,\,\sigma_1^2)$
  & $\mathcal{N}(\mu_2,\,\sigma_2^2)$
  & $\mu_1 \le \mu_2$ and $\sigma_1 = \sigma_2$ \\[4pt]
Inv.\ Gaussian
  & $\mathrm{IG}(\mu_1,\,\lambda_1)$
  & $\mathrm{IG}(\mu_2,\,\lambda_2)$
  & $\lambda_2 \ge \lambda_1$ and $\lambda_1/\mu_1^2 \ge \lambda_2/\mu_2^2$ \\[4pt]
Comp.\ Poisson
  & $\mathrm{CP}(\nu_1,\,F_{Y_1})$
  & $\mathrm{CP}(\nu_2,\,F_{Y_2})$
  & --- (can fail even for $F_{Y_1}{=}F_{Y_2}$, $\nu_1 \le \nu_2$; holds, e.g., for deterministic jumps; see text) \\
\bottomrule
\end{tabular}
\end{table}

\begin{table}[htbp]
\centering
\caption{Conditions for $X \le_{\mathrm{hr}} Y$ (hazard rate order) within each family.}
\label{tab:ordering_hr}
\small
\begin{tabular}{@{}p{2.1cm}p{3.8cm}p{4.0cm}p{4.0cm}@{}}
\toprule
\textbf{Family} & $X$ & $Y$ & \textbf{Condition} \\
\midrule
Gamma
  & $\mathrm{Ga}(\alpha_1,\,\beta_1)$
  & $\mathrm{Ga}(\alpha_2,\,\beta_2)$
  & $\alpha_1 \le \alpha_2$ and $\beta_1 \ge \beta_2$ \\[4pt]
Gaussian
  & $\mathcal{N}(\mu_1,\,\sigma_1^2)$
  & $\mathcal{N}(\mu_2,\,\sigma_2^2)$
  & $\mu_1 \le \mu_2$ and $\sigma_1 = \sigma_2$ \\[4pt]
Inv.\ Gaussian
  & $\mathrm{IG}(\mu_1,\,\lambda_1)$
  & $\mathrm{IG}(\mu_2,\,\lambda_2)$
  & --- (suf.: $\le_{\mathrm{lr}}$ conditions) \\[4pt]
Comp.\ Poisson
  & $\mathrm{CP}(\nu_1,\,F_{Y_1})$
  & $\mathrm{CP}(\nu_2,\,F_{Y_2})$
  & --- (no general criterion; the $\le_{\mathrm{lr}}$ route is unavailable, see text) \\
\bottomrule
\end{tabular}
\end{table}

\begin{table}[htbp]
\centering
\caption{Conditions for $X \le_{\mathrm{st}} Y$ (usual stochastic order) within each family.}
\label{tab:ordering_st}
\small
\begin{tabular}{@{}p{2.1cm}p{3.8cm}p{4.0cm}p{4.0cm}@{}}
\toprule
\textbf{Family} & $X$ & $Y$ & \textbf{Condition} \\
\midrule
Gamma
  & $\mathrm{Ga}(\alpha_1,\,\beta_1)$
  & $\mathrm{Ga}(\alpha_2,\,\beta_2)$
  & $\alpha_1 \le \alpha_2$ and $\beta_1 \ge \beta_2$ \\[4pt]
Gaussian
  & $\mathcal{N}(\mu_1,\,\sigma_1^2)$
  & $\mathcal{N}(\mu_2,\,\sigma_2^2)$
  & $\mu_1 \le \mu_2$ and $\sigma_1 = \sigma_2$ \\[4pt]
Inv.\ Gaussian
  & $\mathrm{IG}(\mu_1,\,\lambda_1)$
  & $\mathrm{IG}(\mu_2,\,\lambda_2)$
  & --- (suf.: $\le_{\mathrm{lr}}$ conditions) \\[4pt]
Comp.\ Poisson
  & $\mathrm{CP}(\nu_1,\,F_{Y_1})$
  & $\mathrm{CP}(\nu_2,\,F_{Y_2})$
  & $\nu_1 \le \nu_2$, $F_{Y_1} = F_{Y_2}$, $Y \ge 0$ \\
\bottomrule
\end{tabular}
\end{table}

\paragraph{Gamma.}
The log-likelihood ratio of $\mathrm{Ga}(\alpha_2,\beta_2)$ over $\mathrm{Ga}(\alpha_1,\beta_1)$ has derivative $(\alpha_2 - \alpha_1)/x - (\beta_2 - \beta_1)$, which is non-negative for all $x > 0$ if and only if $\alpha_2 \ge \alpha_1$ and $\beta_2 \le \beta_1$.
Conversely, if $\alpha_1 > \alpha_2$ then $F_X(t) \sim (\beta_1 t)^{\alpha_1}/(\alpha_1\,\Gamma(\alpha_1)) \ll (\beta_2 t)^{\alpha_2}/(\alpha_2\,\Gamma(\alpha_2)) \sim F_Y(t)$ as $t \to 0^+$, violating $\le_{\mathrm{st}}$;
if $\beta_1 < \beta_2$ then $\bar F_X(t)/\bar F_Y(t) \to \infty$ as $t \to \infty$, again violating $\le_{\mathrm{st}}$.
Hence all three orders coincide for the Gamma family~\cite[Sec.~1.C]{ShakedShanthikumar2007}.
Under the closure condition (same~$\beta$), they all reduce to $\alpha_1 \le \alpha_2$.

\paragraph{Gaussian.}
The CDF of $\mathcal{N}(\mu,\sigma^2)$ is $\Phi\bigl((t-\mu)/\sigma\bigr)$.
The inequality $\Phi\bigl((t-\mu_1)/\sigma_1\bigr) \ge \Phi\bigl((t-\mu_2)/\sigma_2\bigr)$ for all~$t$ requires $(t-\mu_1)/\sigma_1 \ge (t-\mu_2)/\sigma_2$ for all~$t$, which is linear in~$t$ and can hold identically only if $\sigma_1 = \sigma_2$, in which case it reduces to $\mu_1 \le \mu_2$.
The likelihood ratio is then $\exp\bigl((\mu_2-\mu_1)\,x/\sigma^2 + \text{const}\bigr)$, non-decreasing in~$x$.
Thus all three orders coincide: $\sigma_1 = \sigma_2$ and $\mu_1 \le \mu_2$.

\paragraph{Inverse Gaussian.}
The log-likelihood ratio of $\mathrm{IG}(\mu_2,\lambda_2)$ over $\mathrm{IG}(\mu_1,\lambda_1)$ has derivative
\[
  \frac{\lambda_1}{2\mu_1^2} - \frac{\lambda_2}{2\mu_2^2} + \frac{\lambda_2 - \lambda_1}{2x^2}.
\]
This is non-negative for all $x > 0$ if and only if $\lambda_2 \ge \lambda_1$ (from $x \to 0^+$) and $\lambda_1/\mu_1^2 \ge \lambda_2/\mu_2^2$ (from $x \to \infty$).
Under the closure condition (same $\mu^2/\lambda$), these simplify to $\mu_1 \le \mu_2$.
For the general two-parameter case, closed-form necessary and sufficient conditions for $\le_{\mathrm{hr}}$ and $\le_{\mathrm{st}}$ do not appear in the standard references; the $\le_{\mathrm{lr}}$ conditions serve as a sufficient criterion.

\paragraph{Compound Poisson.}
With the same jump-size distribution~$F_Y$ ($Y \ge 0$), $\mathrm{CP}(\nu_1, F_Y) \le_{\mathrm{st}} \mathrm{CP}(\nu_2, F_Y)$ if and only if $\nu_1 \le \nu_2$ (provided $\Pr(Y>0) > 0$): sufficiency follows from the superposition coupling $\mathrm{CP}(\nu_2, F_Y) \stackrel{d}{=} \mathrm{CP}(\nu_1, F_Y) + \mathrm{CP}(\nu_2 - \nu_1, F_Y)$ with independent non-negative summands, and necessity from the atom at zero, $\Pr(X = 0) = e^{-\nu \Pr(Y > 0)}$.
The stronger orders do \emph{not} follow in general: with $Y$ uniform on $\{1, 10\}$ and $\nu_1 = 0.1$, $\nu_2 = 0.2$, the likelihood ratio drops from ${\approx}463$ at $x = 9$ (reachable only by nine unit jumps) to ${\approx}1.8$ at $x = 10$ (reachable by a single large jump), so $\le_{\mathrm{lr}}$ fails, and the survival ratio is likewise non-monotone, so $\le_{\mathrm{hr}}$ fails as well; both hold in special cases such as deterministic jump sizes, where the process is a scaled Poisson.
When the jump-size distributions differ, the ordering conditions depend on the specific form of~$F_Y$ and no universal closed-form characterisation is available.

\paragraph{Implications for the fleet problem.}
Under the closure conditions of Table~\ref{tab:discrete_models}, the cumulative damage $D_{i\ell k}$ depends on the schedule only through cumulative parameters.
For the one-parameter families the comparison of two damage states is scalar:
for Gamma (shared rate~$\beta$), the total shape $\sum_t \alpha_t$;
for IG (shared dispersion $\mu^2/\lambda$), the total mean $\sum_t \mu_t$.
For the Gaussian family the state carries \emph{two} cumulative parameters, $\bigl(\sum_t \mu_t,\ \sum_t \sigma_t^2\bigr)$, and the variance is not schedule-independent: it accumulates with the mission days served and is reduced only by maintenance.
Since the usual stochastic order between Gaussians requires equal variances (Table~\ref{tab:ordering_st}), the Gaussian comparison proceeds through the single-threshold condition~\eqref{eq:loop_gauss} or through componentwise domination of mean and variance, not through the mean alone.
This parameter-level comparison---scalar for Gamma and IG, two-dimensional for Gaussian---is what makes the loop constraint~\eqref{eq:loop} tractable within each model family.

\subsubsection{Loop constraint for each degradation model}

We now specialise the loop constraint~\eqref{eq:loop} to each family.
Let $D_H$ and $D_{2H}$ denote the degradation of a single component at the mid-point and end of the horizon, with parameters determined by the schedule and any intervening maintenance.
The constraint reads $\Pr(D_{2H} > \tau) \le \Pr(D_H > \tau)$.
We drop the $(i,\ell)$ subscripts for clarity.

\paragraph{Gamma.}
Let $D_H \sim \mathrm{Ga}(A_H,\,\beta)$ and $D_{2H} \sim \mathrm{Ga}(A_{2H},\,\beta)$, where $A_k$ denotes the effective cumulative shape at step~$k$ (accounting for repairs and replacements).
The survival function is $\Pr(D > \tau) = Q(A,\,\beta\tau)$, where $Q(a,x) = \Gamma(a,x)/\Gamma(a)$ is the regularised upper incomplete gamma function.
Since $Q(a,x)$ is strictly increasing in~$a$ for fixed $x > 0$, the loop constraint becomes a scalar inequality:
\begin{equation}\label{eq:loop_gamma}
  A_{2H} \;\le\; A_H.
\end{equation}
This is a linear constraint on the cumulative shape parameters, directly expressible in terms of the scheduling and maintenance decision variables.
A caveat on maintenance: the \emph{pathwise} ARD repair $D \mapsto (1-\rho)D$ leaves the shared-rate sub-family, since $(1-\rho)\,\mathrm{Ga}(A,\beta) = \mathrm{Ga}\bigl(A,\,\beta/(1-\rho)\bigr)$.
For the scalar reduction~\eqref{eq:loop_gamma} to remain exact, imperfect repair must be modelled in parameter space, as the shape reduction $A \mapsto (1-\rho)A$ at fixed rate~$\beta$---which matches ARD in expectation, $\mathbb{E}[D^+] = (1-\rho)\,\mathbb{E}[D^-]$, but is a distinct operator pathwise.
The same applies to the inverse Gaussian model below: $c\,\mathrm{IG}(\mu,\lambda) = \mathrm{IG}(c\mu,\,c\lambda)$ multiplies the dispersion $\mu^2/\lambda$ by~$c$, so repair must act as $\mu \mapsto (1-\rho)\mu$ with $\lambda$ induced by the shared dispersion.
The Gaussian family is closed under both readings, $(1-\rho)\,\mathcal{N}(\mu,\sigma^2) = \mathcal{N}\bigl((1-\rho)\mu,\,(1-\rho)^2\sigma^2\bigr)$.

\paragraph{Gaussian.}
Let $D_H \sim \mathcal{N}(\mu_H,\,\sigma_H^2)$ and $D_{2H} \sim \mathcal{N}(\mu_{2H},\,\sigma_{2H}^2)$.
The survival function is $\Pr(D > \tau) = \Phi\bigl((\mu - \tau)/\sigma\bigr)$, where $\Phi$ is the standard normal CDF.
Since $\Phi$ is strictly increasing, the loop constraint is equivalent to
\begin{equation}\label{eq:loop_gauss}
  \frac{\mu_{2H} - \tau}{\sigma_{2H}} \;\le\; \frac{\mu_H - \tau}{\sigma_H},
\end{equation}
or equivalently $(\tau - \mu_{2H})/\sigma_{2H} \ge (\tau - \mu_H)/\sigma_H$.
Note that $\sigma_{2H} = \sigma_H$ does \emph{not} follow from serving the same missions in both halves: absent maintenance the variance accumulates, $\sigma^2_{2H} = \sigma^2_H + \sum_{k = H+1}^{2H} \sigma^2_k > \sigma^2_H$; equality requires the second half's repairs or replacements (which scale the variance by $(1-\rho)^2$ or reset it) to remove exactly the variance the mission days add, in which case~\eqref{eq:loop_gauss} reduces to $\mu_{2H} \le \mu_H$.
In general, condition~\eqref{eq:loop_gauss} couples the mean and standard deviation: a higher variance at~$2H$ can be compensated by a sufficiently lower mean.
Because~\eqref{eq:loop_gauss} constrains only the single scalar $(\mu - \tau)/\sigma$ of a two-parameter state, it does not by itself make the schedule repeatable (cf.\ Section~\ref{sec:loop_routeI}); repeatability is restored by the componentwise condition $\mu_{2H} \le \mu_H$ and $\sigma^2_{2H} \le \sigma^2_H$---the exact-law analogue of~\eqref{eq:loop_moments}---which implies~\eqref{eq:loop_gauss} whenever $\mu_H < \tau$ and chains because both moments propagate monotonically under the dynamics.

\paragraph{Inverse Gaussian.}
Let $D_H \sim \mathrm{IG}(\mu_H,\,\lambda_H)$ and $D_{2H} \sim \mathrm{IG}(\mu_{2H},\,\lambda_{2H})$.
The CDF of $\mathrm{IG}(\mu,\lambda)$ admits the closed form
\begin{equation}\label{eq:ig_cdf}
  F_{\mathrm{IG}}(x;\,\mu,\lambda) \;=\;
    \Phi\!\Biggl(\sqrt{\frac{\lambda}{x}}\Bigl(\frac{x}{\mu}-1\Bigr)\Biggr)
    \;+\; e^{2\lambda/\mu}\;
    \Phi\!\Biggl(-\sqrt{\frac{\lambda}{x}}\Bigl(\frac{x}{\mu}+1\Bigr)\Biggr).
\end{equation}
The loop constraint becomes
\begin{equation}\label{eq:loop_ig}
  F_{\mathrm{IG}}(\tau;\,\mu_{2H},\lambda_{2H}) \;\ge\; F_{\mathrm{IG}}(\tau;\,\mu_H,\lambda_H).
\end{equation}
This is an analytical but nonlinear inequality in the four parameters $(\mu_H,\lambda_H,\mu_{2H},\lambda_{2H})$.
Under the closure condition (same $\mu^2/\lambda$), the $\le_{\mathrm{lr}}$ results of Table~\ref{tab:ordering_lr} guarantee that $\mu_{2H} \le \mu_H$ is a sufficient condition.

\paragraph{Compound Poisson.}
When $D \sim \mathrm{CP}(\nu,F_Y)$, the survival function does not admit a universal closed form.
Conditioning on the number of shocks gives the series
\begin{equation}\label{eq:surv_cp}
  \Pr(D > \tau) \;=\; \sum_{n=0}^{\infty} \frac{e^{-\nu}\,\nu^n}{n!}\;\Pr\!\Bigl(\sum_{k=1}^{n} Y_k > \tau\Bigr),
\end{equation}
which can be evaluated numerically (it converges rapidly) or bounded.

\paragraph{Upper bounds when no closed form is available.}
When the survival function $\Pr(D > \tau)$ cannot be computed in closed form---as for the compound Poisson process, or when the closure condition fails and the marginal of the sum is not within a standard family---the \emph{reliability} constraint can still be enforced through conservative upper bounds on the failure probability: Markov ($\Pr(D > \tau) \le \mathbb{E}[D]/\tau$), Cantelli~\eqref{eq:cantelli_surrogate}, and Chernoff~\eqref{eq:chernoff_surrogate}.
Section~\ref{sec:cons_surrogates} develops these bounds in per-mission form, together with the Bernstein and Hoeffding refinements, their relative tightness in the small-$\varepsilon$ regime, and their MILP encodings; for the compound Poisson process the Chernoff bound is particularly convenient, since $M_D(s) = \exp\bigl(\nu\,(M_Y(s)-1)\bigr)$ is available in closed form whenever $M_Y$ is.
The \emph{loop} constraint, by contrast, must not be translated into the same bounds: imposing $U_{2H} \le U_H$ on a scalar bound~$U$ does not make the schedule repeatable, and the loop is instead imposed on the componentwise moment condition~\eqref{eq:loop_moments}, whose two sides are additive over increments and therefore linear in the scheduling decisions (Section~\ref{sec:loop_routeI}).

\subsection{Maintenance strategies}
\label{sec:maint_strategies}

Maintenance actions on a repairable component range between two extremes: \emph{minimal repair}, which restores the component to the state it was in just before failure (``as bad as old'', ABAO), and \emph{perfect repair} or \emph{replacement}, which resets it to a like-new condition (``as good as new'', AGAN).
Real-world interventions almost always fall between these two extremes; they are collectively referred to as \emph{imperfect maintenance}~\cite{PhamWang1996}.
This section surveys the principal modelling families.

\subsubsection{Taxonomy of repair actions}

Table~\ref{tab:repair_actions} classifies five canonical maintenance actions by their effect on the component state.

\begin{table}[htbp]
\centering
\caption{Taxonomy of maintenance actions, from least to most restorative.}
\label{tab:repair_actions}
\small
\begin{tabular}{@{}p{3.0cm}p{4.5cm}p{5.8cm}@{}}
\toprule
\textbf{Action} & \textbf{Effect on damage $D$} & \textbf{Model} \\
\midrule
No intervention
  & $D^+ = D^-$
  & Damage unchanged \\[3pt]
Minimal repair (ABAO)
  & $D^+ = D^-$
  & Failure intensity restored to pre-failure level; damage history preserved~\cite{BrownProschan1983} \\[3pt]
Imperfect repair
  & $D^- > D^+ > D^{\mathrm{new}}$
  & Partial reduction of damage or virtual age (see below) \\[3pt]
Perfect repair (AGAN)
  & $D^+ = D^{\mathrm{new}}$
  & Component returned to like-new state \\[3pt]
Replacement
  & $D^+ = D^{\mathrm{new}}$
  & Physical component substituted; equivalent to perfect repair \\
\bottomrule
\end{tabular}
\end{table}

\noindent Here $D^-$ and $D^+$ denote the damage immediately before and after maintenance, and $D^{\mathrm{new}}$ is the damage of a new component.

\subsubsection{Virtual age models}

Kijima~\cite{Kijima1989} formalised imperfect repair through the concept of \emph{virtual age}.
After the $n$-th repair, the component behaves as if it had age $V_n$ (rather than its calendar age), and the failure intensity at calendar time~$t$ is $\lambda(t - T_n + V_n)$, where $T_n$ is the time of the $n$-th repair and $\lambda(\cdot)$ is the baseline intensity.
Two variants are standard:
\begin{itemize}
  \item \textbf{Kijima~I.} The repair acts only on the damage accumulated since the previous repair:
    \[
      V_n \;=\; V_{n-1} + (1-\rho)\,(T_n - T_{n-1}),
    \]
    where $\rho \in [0,1]$ is the repair efficiency.
    Only the most recent inter-repair period is partially rejuvenated (``memory-one'').
  \item \textbf{Kijima~II.} The repair acts on the entire accumulated virtual age:
    \[
      V_n \;=\; (1-\rho)\,(V_{n-1} + T_n - T_{n-1}).
    \]
    The full history is rejuvenated (``infinite memory'').
\end{itemize}
In both cases $\rho = 0$ gives minimal repair (ABAO) and $\rho = 1$ gives perfect repair (AGAN).

\subsubsection{Reduction-based models}

Doyen and Gaudoin~\cite{DoyenGaudoin2004} proposed a unifying taxonomy by distinguishing the \emph{quantity} that the repair reduces: the accumulated damage (or virtual age), or the failure intensity.
Each class admits a memory-one variant (subscript~1) and an infinite-memory variant (subscript~$\infty$).

\paragraph{Arithmetic Reduction of Damage / Age (ARD, ARA).}
The repair removes a fraction~$\rho$ of the accumulated damage (ARD) or virtual age (ARA):
\begin{align}
  \text{ARD}_1: &\quad D^+ = D^- - \rho\,(D^- - D^{\mathrm{last}}), \label{eq:ard1} \\
  \text{ARA}_1: &\quad V^+ = V^- - \rho\,(V^- - V^{\mathrm{last}}), \label{eq:ara1} \\
  \text{ARD}_\infty: &\quad D^+ = (1-\rho)\,D^-, \label{eq:ardinf} \\
  \text{ARA}_\infty: &\quad V^+ = (1-\rho)\,V^-, \label{eq:arainf}
\end{align}
where $D^{\mathrm{last}}$ (resp.~$V^{\mathrm{last}}$) is the damage (resp.\ virtual age) at the end of the previous maintenance epoch.
$\text{ARD}_\infty$ coincides with Kijima~II expressed in damage units, and $\text{ARA}_1$ coincides with Kijima~I expressed in virtual age.

\paragraph{Arithmetic Reduction of Intensity (ARI).}
Instead of reducing the virtual age, the repair directly reduces the failure intensity:
\begin{align}
  \text{ARI}_1: &\quad \lambda^+ = \lambda^- - \rho\,(\lambda^- - \lambda^{\mathrm{last}}), \label{eq:ari1} \\
  \text{ARI}_\infty: &\quad \lambda^+ = (1-\rho)\,\lambda^-. \label{eq:ariinf}
\end{align}
ARI models are natural when maintenance targets the failure rate (e.g.\ condition-based interventions) rather than the underlying degradation state.
They are in general not equivalent to ARA/ARD models: a reduction in intensity does not correspond to a unique reduction in virtual age unless the baseline intensity is strictly monotone~\cite{DoyenGaudoin2004}.

\subsubsection{Connection to the fleet management formulation}

The abstract repair operator $\mathcal{R}(D^-,\,\rho)$ introduced in Section~\ref{sec:overview} encompasses ARD and replacement models.
In particular:
\begin{itemize}
  \item $\text{ARD}_\infty$: $\mathcal{R}(D^-,\rho) = (1-\rho)\,D^-$ (used in the scoping as the default ``ARD'' policy);
  \item $\text{ARD}_1$: $\mathcal{R}(D^-,\rho) = D^- - \rho\,(D^- - D^{\mathrm{last}})$, requiring the additional state variable $D^{\mathrm{last}}$;
  \item Replacement / AGAN: $\mathcal{R}(D^-,\rho) = D^{\mathrm{new}}$ (captured by the separate replacement variable $x^r_{i\ell k}$ in the formulation).
\end{itemize}
The choice of repair model affects the degradation dynamics~\eqref{eq:dyn_full}, the definition of the repair cost variable~$z_{i\ell k}$, and the tractability of the resulting optimisation problem.
ARI models are not directly compatible with the damage-tracking formulation of Section~\ref{sec:overview} and would require reformulation in terms of the failure intensity rather than the damage state.

\paragraph{Maintenance under the Palmgren--Miner damage model.}
The remaining-life framework of Section~\ref{sec:discrete_deg} represents damage as the normalised fraction $D_{\mathrm{PM}} = \sum_i n_i/N_{f,i} \in [0,1]$, with failure at $D_{\mathrm{PM}} \ge 1$.
Because $D_{\mathrm{PM}}$ is additive over loading blocks, it fits the same damage-tracking structure as the stochastic processes of Table~\ref{tab:deg_models}: each mission day contributes an increment $\Delta D_{\mathrm{PM}} = n_j / N_{f,j}$ (deterministic if the S--N parameters are fixed, or random if scatter is modelled~\cite{MunizCalvente2022}).
The repair operator~$\mathcal{R}$ and the replacement action therefore apply directly:
\begin{itemize}
  \item \emph{Replacement} resets the fatigue fraction to zero: $D_{\mathrm{PM}}^+ = 0$ (equivalently $D^{\mathrm{new}} = 0$).
    This corresponds to substituting the physical component with a new one whose cycle count is zero.
  \item $\text{ARD}_\infty$: $D_{\mathrm{PM}}^+ = (1-\rho)\,D_{\mathrm{PM}}^-$.
    Physically, this models an intervention that partially restores fatigue life (e.g.\ surface treatment, shot peening, localised crack repair) without resetting the full loading history.
  \item $\text{ARD}_1$: $D_{\mathrm{PM}}^+ = D_{\mathrm{PM}}^- - \rho\,(D_{\mathrm{PM}}^- - D_{\mathrm{PM}}^{\mathrm{last}})$.
    Only the damage accumulated since the last maintenance epoch is partially reversed.
\end{itemize}
A notable difference from the stochastic process models is that $D_{\mathrm{PM}}$ is bounded above by construction (failure at~1), so the threshold is fixed at $\tau = 1$ and the reliability constraint becomes $\Pr(D_{\mathrm{PM}} \ge 1) \le \varepsilon$.
In the deterministic variant of Miner's rule this is a hard feasibility check; in the stochastic extension it becomes a chance constraint analogous to~\eqref{eq:abs_rel}.

Closure under summation is trivially satisfied: $D_{\mathrm{PM}}$ is a sum of real-valued increments, so the marginal at any step is the convolution of the individual increments regardless of their distributional family.
In the deterministic variant of Miner's rule, comparing two schedules reduces to comparing two real numbers and no stochastic ordering is needed.
In the stochastic extension (scatter in $N_{f,i}$), the parametric ordering conditions of Section~\ref{sec:stoch_orders} apply only if the distributional assumption on the scatter places the increments $n_i/N_{f,i}$ in one of the families of Table~\ref{tab:deg_models}; otherwise the general ordering definitions still hold but must be verified case by case.

\paragraph{Remaining items for this section.}
\begin{itemize}
    \item[$\checkmark$] Fatigue models (either damage accumulation or remaining life)
    \item[$\square$] Degradation-aware fleet management
    \item[$\checkmark$] Maintenance strategies
    \item[$\checkmark$] Stochastic orderings
    \item[$\square$] Similar applications
    \item[$\square$] Novelties of the project and its placement in the literature
\end{itemize}

\section{Focus on electric vehicles}

A battery-electric road vehicle---a transit bus or a passenger car---is a complex repairable system whose parts degrade at very different rates and through very different physical mechanisms.
Not every part deserves an explicit place in the fleet-management model of Section~\ref{sec:overview}.
We retain a component if it satisfies two criteria:
\begin{enumerate}[label=(C\arabic*)]
  \item \textbf{Usage-driven degradation.}
    Its condition deteriorates measurably as a function of \emph{how} the vehicle is used, so that the mission assigned on a given day produces a quantifiable damage increment $\Delta D^{(j)}_{i\ell k}$ (Section~\ref{sec:overview}).
    A part that ages only with calendar time, independently of the mission served, carries no scheduling information and need not enter the per-mission degradation state.
  \item \textbf{Operational significance.}
    Its failure or preventive renewal is a material driver of at least one of \emph{safety}, \emph{availability}, or \emph{life-cycle maintenance cost}---the three quantities the schedule trades off.
\end{enumerate}
We deliberately do \emph{not} require that the component admit a degradation model expressible within the families of Section~\ref{sec:stoch_orders}: a component can be worth considering even when the abstract model of Section~\ref{sec:overview} does not yet capture it.
Such cases are listed here together with the modelling extension they would demand, rather than being silently dropped.

The components meeting these criteria fall into two broad groups.
The \emph{electrochemical and electrical powertrain}---traction battery, electric motor, power electronics, charging interface---degrades through electrochemical ageing and thermal and electrical cycling, driven by energy throughput, charge and discharge rates, and temperature.
The \emph{mechanical running gear}---tires, friction brakes, wheel bearings, suspension, reduction gear---degrades through wear and fatigue, driven by distance travelled, vehicle mass, road quality, and the acceleration/braking profile.
Table~\ref{tab:ev_components} lists the principal candidates in each group, the dominant degradation mechanism, the operational quantity that drives it, and the reason the component is operationally significant.

Two domain-specific facts shape the relative importance of these components.
First, the traction battery dominates: a recent predictive-maintenance study of electric buses centres entirely on the battery pack as the primary maintenance-critical component, models its end-of-life through the accumulated charge throughput, and shows that optimising the driving profile alone can cut the modelled capacity loss by 25\%~\cite[Secs.~2,~4.2]{Ibrahim2025}---direct evidence that battery damage is steerable by exactly the kind of mission-assignment decisions our framework optimises.
Second, regenerative braking absorbs most of the service braking, with estimated reductions of roughly 60--90\% in friction-brake wear emissions~\cite[Sec.~6.1]{Fussell2022}; the relative wear burden consequently shifts to the tires, whose wear is aggravated by the higher mass of electric vehicles---to the point that, on motorways, the added tire-wear particle emissions are not offset by the regenerative-braking savings~\cite[Sec.~6.1]{Fussell2022}.

\begin{table}[htbp]
\centering
\caption{Candidate electric-vehicle components for degradation-aware modelling.
Relevance flags: \textbf{S} safety-critical, \textbf{A} availability driver (its failure or renewal takes the vehicle out of service), \textbf{C} maintenance-cost driver.
The ``mission driver'' column names the operational quantity that determines the per-mission damage increment $\Delta D^{(j)}_{i\ell k}$.}
\label{tab:ev_components}
\small
\begin{tabular}{@{}p{3.0cm}p{4.6cm}p{3.5cm}c@{}}
\toprule
\textbf{Component} & \textbf{Dominant degradation mechanism} & \textbf{Mission driver} & \textbf{Rel.} \\
\midrule
\multicolumn{4}{@{}l}{\textit{Electrochemical / electrical powertrain}} \\
Traction battery pack
  & Capacity fade and impedance growth: SEI growth, lithium plating, loss of active material (cycle $+$ calendar ageing)
  & Energy throughput, DOD, C-rate, temperature
  & A, C \\[3pt]
Traction motor
  & Winding-insulation thermal ageing; bearing wear
  & Thermal/load cycles, current
  & A, C \\[3pt]
Power semiconductors (IGBT / SiC modules in the inverter and converters)
  & Thermal-cycling fatigue (bond-wire lift-off, solder cracking)
  & Power throughput, thermal cycles
  & A, C \\[3pt]
Converter capacitors (DC-link, filtering)
  & Electrolyte vaporisation / dielectric wear-out (ripple-current self-heating)
  & Ripple current, temperature, operating hours
  & A, C \\[3pt]
Control / power PCBs (gate drivers, BMS, converter boards)
  & Solder-joint and via thermo-mechanical fatigue; random infant/wear-out failures
  & Thermal cycles, power-on hours, vibration
  & A \\[3pt]
Charging interface (inlet, connector)
  & Contact wear and contact-resistance growth over mating cycles
  & Number of charging events
  & A \\
\midrule
\multicolumn{4}{@{}l}{\textit{Mechanical: running gear and chassis}} \\
Tires
  & Tread abrasion (friction work at the contact patch)
  & Distance, vehicle mass, driving style, road surface
  & S, C \\[3pt]
Brake friction elements (pads, discs)
  & Abrasive and thermal wear (mitigated by regenerative braking)
  & Number and energy of friction-braking events
  & S, C \\[3pt]
Wheel bearings
  & Rolling-contact fatigue spalling
  & Distance, load
  & S, A \\[3pt]
Suspension (springs, dampers, bushings)
  & Fatigue; loss of damping
  & Load cycles, road quality
  & A \\[3pt]
Reduction gear / drive shafts
  & Wear, contact fatigue, lubricant ageing
  & Transmitted torque, distance
  & A \\[3pt]
Auxiliaries (HVAC, battery thermal management, doors)
  & Wear, cycle-driven failures
  & Operating hours, actuation cycles
  & A \\
\bottomrule
\end{tabular}
\end{table}

\paragraph{Power-electronic and electronic components.}
The traction chain contains four electrical sub-families beyond the battery, and they fit the framework to differing degrees.
The \emph{power semiconductors} of the inverter and the DC--DC and on-board-charger \emph{converters} are governed by junction-temperature-swing fatigue and map cleanly onto the remaining-life family (treated in detail below).
The \emph{converter capacitors}---DC-link and filtering---wear out by electrolyte vaporisation and dielectric ageing driven by ripple-current self-heating; this is a monotone, usage-driven accumulation (Arrhenius in temperature, power law in ripple current) and therefore also fits the framework, with the mission acting through the delivered power and the resulting ripple and thermal load.\footnote{The capacitor and PCB degradation-model forms in this paragraph are stated from standard power-electronics reliability theory; located primary sources have not yet been secured for them (see \texttt{inaccessible\_papers.md}).}
The \emph{control and power printed circuit boards}---gate-driver, battery-management, and converter boards---are the borderline case: their solder-joint and plated-through-via fatigue is genuinely usage-driven (thermal cycling, power-on hours, vibration) and so satisfies~(C1)--(C2), but a substantial part of their failures are sudden and effectively random rather than a smooth accumulation crossing a threshold.
The damage-accumulation dynamics of Section~\ref{sec:overview} do not fully capture this mixed wear-out-plus-random behaviour; representing it faithfully would require either a compound-Poisson shock term (Section~\ref{sec:stoch_orders}) or a failure-intensity formulation rather than the damage-state formulation of Equation~\eqref{eq:dyn_full}.
We list PCBs here, with this caveat, precisely so the limitation is recorded rather than hidden by the choice of model.

\paragraph{Modelling shortlist.}
For a first concrete instantiation of the abstract framework, the most informative components are those that are simultaneously safety- or cost-critical \emph{and} equipped with mature, monotone degradation models that fit the closure and ordering analysis of Section~\ref{sec:stoch_orders}: the \emph{traction battery pack}, the \emph{tires}, the \emph{brake friction elements}, the \emph{wheel bearings}, the \emph{traction-motor insulation}, and the \emph{inverter and converter power modules}.
Each maps onto a single per-component degradation variable $D_{i\ell k}$ whose mission-dependent increment $\Delta D^{(j)}_{i\ell k}$ is governed by the operational driver in the third column of Table~\ref{tab:ev_components}.
A long highway mission, for instance, produces large tire-wear and deep-discharge battery increments, whereas a stop-intensive urban mission multiplies braking events and charge--discharge micro-cycles; this mission-dependence is exactly the coupling that the scheduling problem of Section~\ref{sec:overview} exploits.

\paragraph{Worked example: traction-battery capacity fade.}
The battery pack is the costliest component of an electric vehicle and the focus of its predictive maintenance~\cite[Secs.~2,~4.2]{Ibrahim2025}, so it serves as the reference instantiation of the abstract increment $\Delta D^{(j)}_{i\ell k}$.
Its degradation state is the relative capacity fade (one minus state-of-health), and the failure threshold~$\tau$ is the maximum admissible fade (the end-of-life criterion).
Mechanistically, capacity is lost through three side-reaction families---SEI film growth, lithium plating, and gas generation---whose contributions split into loss of lithium inventory and loss of active material~\cite[Secs.~2.1,~4.1, Eqs.~28--29]{Lv2025}.
At the fleet level this is abstracted by semi-empirical models that decompose the fade into a \emph{calendar} term and a \emph{cycling} term: the calendar term grows with an Arrhenius factor $a\exp(-E_a/RT)$ in temperature and a power law in time and state-of-charge, while the cycling term grows with the number of cycles modulated by temperature, C-rate, and depth-of-discharge stress factors~\cite[Eqs.~1--3]{Hosen2021}.
An even simpler engineering abstraction links end-of-life directly to the accumulated charge throughput~\cite[Sec.~4.2]{Ibrahim2025}.
For our framework the decisive structure is: the cycling fade is monotone and additive over days, and the daily increment is a function of the energy served on the mission---including the recharge that restores it, so that fast charging (high C-rate) makes the same mission costlier in damage---together with its depth-of-discharge and temperature profile.
This is precisely a mission-dependent accumulation process of the Gamma or inverse-Gaussian type of Section~\ref{sec:stoch_orders}; the non-stationary (power-law) time dependence is absorbed by the time-varying increment parameters that the formulation already allows ($\Delta D^{(j)}_{i\ell k}$ may depend on~$k$).
The calendar term, by contrast, accrues even on idle and maintenance days and is therefore \emph{not} captured by the dynamics of Section~\ref{sec:overview}, which leave the state unchanged on such days; over short horizons it can be neglected, while over long horizons the dynamics~\eqref{eq:dyn_full} must be extended with a small activity-independent drift added to every case, including idle.
Maintenance also takes a specific form: replacing a subset of battery modules removes a corresponding fraction of the pack-level damage---a physical realisation of the ARD repair operator $\mathcal{R}$ with $\rho$ the fraction of modules renewed---while full pack replacement is the AGAN action~$x^r$.

\subsection{Degradation models for the shortlisted components}
\label{sec:ev_deg_models}

The remaining shortlisted components degrade through the same two structural mechanisms identified for the battery---monotone usage-driven \emph{wear/accumulation} and \emph{cyclic-load fatigue life consumption}---and therefore map onto the same two model families of Section~\ref{sec:stoch_orders}.
Table~\ref{tab:ev_deg_models} summarises the governing model for each, the operational quantity that drives the per-mission increment, and the abstract family it instantiates.

\begin{table}[htbp]
\centering
\caption{Degradation models for the shortlisted electric-vehicle components and their mapping onto the abstract model families of Section~\ref{sec:stoch_orders}.
``Accumulation'' denotes a monotone damage-accumulation process (Gamma / inverse-Gaussian / compound-Poisson); ``remaining-life'' denotes a cycles-to-failure model accumulated by the Palmgren--Miner rule.}
\label{tab:ev_deg_models}
\small
\begin{tabular}{@{}p{2.4cm}p{4.9cm}p{3.0cm}p{2.6cm}@{}}
\toprule
\textbf{Component} & \textbf{Degradation model / law} & \textbf{Mission driver} & \textbf{Abstract family} \\
\midrule
Battery pack
  & Calendar $+$ cycle capacity fade: Arrhenius in $T$, power law in $t$/SOC, cycle-count term \cite[Eqs.~1--3]{Hosen2021}; mechanistic SEI / plating / LAM \cite[Sec.~2.1]{Lv2025}
  & Energy throughput, DOD, C-rate, temperature
  & Accumulation ($k$-dependent) \\[3pt]
Tires
  & Friction-work abrasion \cite[Eq.~7]{SchallamachTurner1960}; factors: tire build, road surface, vehicle operation \cite[Sec.~2.1.2]{Fussell2022}
  & Distance, mass, accel./braking, road surface
  & Accumulation \\[3pt]
Brake friction pad
  & Archard-type wear per braking event \cite[Eqs.~14,~22]{Sha2022}; regenerative-braking mitigation \cite[Sec.~6.1]{Fussell2022}
  & Number and energy of friction-braking events
  & Accumulation \\[3pt]
Wheel bearing
  & Subsurface-fatigue spall propagation, continuum damage \cite[Sec.~2.1.3]{Ohana2023}
  & Distance, speed, load
  & Accumulation / remaining-life \\[3pt]
Traction-motor insulation
  & Arrhenius--Dakin thermal-endurance life \cite{Dakin1948}; Montsinger rule \cite{Montsinger1930}; cumulative-damage extension \cite[Eqs.~1--2]{Zhang2025insul}
  & Thermal / load cycles, current
  & Remaining-life (Miner) \\[3pt]
Inverter power module
  & Coffin--Manson / Norris--Landzberg cycles-to-failure in $\Delta T_j$ \cite[Eqs.~3--5]{Alavi2024}; Bayerer/CIPS08 \cite[Eq.~1]{Bayerer2008}
  & Power throughput, thermal cycles
  & Remaining-life (Miner) \\
\bottomrule
\end{tabular}
\end{table}

\paragraph{Tires.}
Tread loss is governed by the frictional work expended at the tire--road contact patch: the classical theory of Schallamach and Turner assumes abrasion proportional to the frictionally dissipated energy~\cite[Eq.~7]{SchallamachTurner1960} and derives that, at small slip, wear grows approximately with the \emph{square} of the slip~\cite[Eq.~19]{SchallamachTurner1960}---so the wear rate rises steeply with the tangential forces generated by traction, braking, and cornering.
The principal factors are the tire characteristics (compound, construction), the road-surface characteristics, and the vehicle operation~\cite[Sec.~2.1.2]{Fussell2022}.
The higher kerb mass of electric vehicles raises these contact forces---enough that, on motorways, the added tire-wear particle emissions are not offset by the brake-wear savings of regenerative braking~\cite[Sec.~6.1]{Fussell2022}.
Tread depth decreases monotonically from its new value to the legal limit, which plays the role of the threshold~$\tau$; the per-mission increment depends on distance, load, road surface, and driving aggressiveness.
For bus fleets, retreading acts as a near-AGAN replacement of the wear state, while tire rotation merely redistributes wear across positions and is left outside the per-component model.

\paragraph{Brake friction elements.}
Friction-pad material loss is governed by the energy dissipated at the pad--disc interface during braking; a physics-based route applies Archard's law on a fractal, elastoplastic contact to obtain a per-cycle wear depth and hence a brake-cycle life $N = H_0/(S_0\,h_{\mathrm{wear}})$, with $H_0$ the admissible wear depth and $h_{\mathrm{wear}}$ the wear per braking~\cite[Eqs.~14,~22]{Sha2022}.
In an electric vehicle the regenerative brake absorbs most service braking, cutting friction-brake wear by an estimated 60--90\%~\cite[Sec.~6.1]{Fussell2022}; the residual increments concentrate in the events where regeneration is unavailable or insufficient---battery near full charge, very low temperature, emergency stops.
The cumulative pad wear remains a sum of non-negative, mission-dependent increments---one per friction-braking event---i.e.\ a monotone accumulation process whose rate now also depends on the energy-management strategy.

\paragraph{Wheel bearings.}
Bearing degradation proceeds in two stages: a long rolling-contact-fatigue \emph{initiation} phase followed by rapid \emph{spall propagation} once a surface defect forms.
Ohana et al.~\cite[Sec.~2.1.3]{Ohana2023} model the propagation phase with a continuum-damage-mechanics formulation in which a scalar damage variable accumulates per over-rolling cycle, $\mathrm{d}D/\mathrm{d}N = \bigl(\sigma_{\mathrm{eff}}/[\sigma_r(1-D)]\bigr)^m$, with growth driven by the tensile residual stresses generated at the spall edge by repeated rolling-element impacts; the model is intended as the basis for a physics-based prognostic tool that estimates remaining useful life once a critical spall size is reached.
Because the underlying driver is the number of load cycles---and hence running distance and speed---the bearing damage is a monotone accumulation process, well represented by a Gamma or inverse-Gaussian process (initiation) feeding a remaining-life threshold (propagation).

\paragraph{Traction-motor insulation.}
The winding insulation is the principal life-limiting element of the traction motor, and its dominant ageing mechanism is thermal: elevated winding temperature accelerates the chemical degradation of the dielectric.
Dakin recognised thermal ageing as a chemical reaction whose rate constant follows the Arrhenius exponential law $K = A\,e^{-B/T}$ in absolute temperature~\cite[Eq.~5]{Dakin1948}, giving the thermal-endurance life $L(T) = a\,\exp(b/T)$~\cite[Eq.~1]{Zhang2025insul}; the earlier empirical rule of Montsinger states that the ageing rate roughly doubles for each 8\,$^\circ$C temperature increase~\cite[Synopsis, p.~776]{Montsinger1930}---a halving of life per fixed rise that Dakin showed to be the constant-coefficient approximation of the Arrhenius law, accurate at one temperature only~\cite[Fig.~1 and discussion of Eq.~5]{Dakin1948}.
Under a mission-varying temperature profile, the fraction of life consumed in an interval $\mathrm{d}t$ at temperature $T_i(t)$ is $\mathrm{d}t/L[T_i(t)]$, and the total consumed fraction is the integral of this loss rate over the profile---a cumulative-damage (Palmgren--Miner) extension of the Arrhenius model~\cite[Eq.~2 and Fig.~10]{Zhang2025insul}.
The abstract description therefore coincides with the remaining-life model of Section~\ref{sec:discrete_deg}, with the mission driving the thermal load (current, cooling, ambient); notably, the same review also surveys Wiener- and Gamma-process formulations of insulation degradation~\cite[Sec.~5]{Zhang2025insul}, directly matching the stochastic families of Section~\ref{sec:stoch_orders}.

\paragraph{Inverter power modules.}
The IGBT or SiC modules of the traction inverter fail predominantly by thermo-mechanical fatigue---aluminium bond-wire lift-off and solder-layer cracking---driven by the coefficient-of-thermal-expansion mismatch under repeated junction-temperature swings $\Delta T_j$~\cite[Sec.~1]{Bayerer2008}.
The number of cycles to failure is classically captured by a Coffin--Manson power law $N_f = \alpha\,(\Delta T_j)^{-n}$~\cite[Eq.~3]{Alavi2024}, refined by adding an Arrhenius factor in the mean junction temperature $T_{jm}$, $N_f = A\,(\Delta T_j)^{\alpha}\exp\!\bigl(E_a/(k_b T_{jm})\bigr)$~\cite[Eq.~4]{Alavi2024}, and generalised by the multi-factor Bayerer (CIPS08) model
$N_f = K\,\Delta T_j^{\beta_1} e^{\beta_2/(T_{j}+273)}\, t_{\mathrm{on}}^{\beta_3}\, I^{\beta_4}\, V^{\beta_5}\, D^{\beta_6}$,
fitted by multiple regression over power-cycling tests spanning temperature swing, power-on time~$t_{\mathrm{on}}$, current per bond wire~$I$, blocking voltage~$V$, and bond-wire diameter~$D$~\cite[Eq.~1]{Bayerer2008}.
Being purely empirical, the model is valid only within the fitted test range and is not to be extrapolated~\cite[Sec.~3.1]{Bayerer2008}; its coefficients also require recalibration per device family---a recent re-fit on dedicated power-cycling tests shifted $(\beta_1,\beta_2,\beta_3)$ substantially and cut the predicted B10 inverter lifetime from 43 to 13~years~\cite[Abstract and Sec.~5]{Alavi2024}.
Under a realistic mission mix the thermal cycles have heterogeneous amplitudes, and the accumulated damage is summed by the Palmgren--Miner rule $Q = \sum_i n_i/N_{f,i}$, failure at $Q=1$---placing the power module squarely in the remaining-life family of Section~\ref{sec:discrete_deg}, with the mission determining the transmitted power and hence the temperature-swing spectrum.

\paragraph{Common structure.}
Every shortlisted component reduces to one of the two abstract descriptions of Section~\ref{sec:stoch_orders}: a \emph{monotone accumulation process} (battery cycling fade, tires, brakes, bearings) or a \emph{remaining-life model accumulated by the Palmgren--Miner rule} (insulation, power modules).
In both, the per-mission damage increment $\Delta D^{(j)}_{i\ell k}$ is set by an operational driver that the schedule controls---energy throughput and charging rate, distance and load, braking profile, or thermal cycling---which is precisely the structure the abstract formulation of Section~\ref{sec:overview} assumes.
This confirms that the electric-vehicle fleet is a faithful concrete instance of the general problem, and fixes which degradation family parameterises each component when the model is specialised in Section~\ref{sec:evs}.

\paragraph{Excluded parts.}
Components whose deterioration is essentially independent of the mission served---body corrosion, interior fittings, static low-voltage electronics, the 12-volt auxiliary battery---satisfy criterion~(C2) but not~(C1).
They are best handled by time-based overhaul outside the optimisation and are therefore omitted from the per-mission degradation state, although their fixed maintenance windows can be imposed as exogenous availability constraints.
The calendar-ageing component of the traction battery is the one borderline case and is handled by the drift extension discussed in the worked example above.

\paragraph{Remaining items for this section.}
\begin{itemize}
    \item[$\checkmark$] Relevant components (mechanical and electrical) to model an electric vehicle
    \item[$\checkmark$] Typical degradation model for each component, and the operational assumptions under which it holds
\end{itemize}

\section{Optimization problem for a fleet of electric vehicles}
\label{sec:evs}

Table~\ref{tab:notation_evs} lists the additional symbols introduced when specialising the abstract formulation of Section~\ref{sec:overview} to a concrete fleet of electric vehicles.
All symbols from Table~\ref{tab:notation} remain in force; the entries below either instantiate abstract quantities or add new ones.

\begin{table}[htbp]
\centering
\caption{Additional notation for the electric-vehicle fleet specialisation (extends Table~\ref{tab:notation}).}
\label{tab:notation_evs}
\small
\begin{tabular}{@{}lp{10.5cm}@{}}
\toprule
\textbf{Symbol} & \textbf{Description} \\
\midrule
\multicolumn{2}{@{}l}{\textit{Cost parameters}} \\
$C_M$                        & Cost of scheduling a maintenance day (scalar) \\
$C_D$                        & Damage regularisation coefficient (scalar) \\
$C_{R,i\ell}$                & Cost of imperfect repair per unit of expected degradation removed \\
$C_{\mathrm{rep},i\ell}$     & Cost of full replacement of component $\ell$ of member $i$ \\
\bottomrule
\end{tabular}
\end{table}

\subsection{Cost function}

With the cost parameters of Table~\ref{tab:notation_evs}, the abstract objective of Section~\ref{sec:overview} takes the form:
\begin{align}\label{eq:objective_evs}
\min \quad
  & C_M \sum_{k=1}^{2H} \sum_{i=1}^{F} y_{ik}
  \;+\; \sum_{k=1}^{2H} \sum_{i=1}^{F} \sum_{\ell=1}^{L} C_{R,i\ell}\, z_{i\ell k}
  \;+\; \sum_{k=1}^{2H} \sum_{i=1}^{F} \sum_{\ell=1}^{L} C_{\mathrm{rep},i\ell}\, x^r_{i\ell k} \notag \\
  & +\; C_D\, u.
\end{align}
Each term is a weighted sum over the planning horizon:
the first counts maintenance days, the second scales with the degradation removed by each repair, the third counts replacements, and the last penalises the worst-case aggregate expected damage via~$u$ (cf.~\eqref{eq:u_def}).

\begin{itemize}
    \item Specialize the previous optimization problem to the case of electric vehicles
    \item Look for what are can be exactly modeled and what no
    \item Provide the MILP approximation of such problem
\end{itemize}

\section{Implementation details}

\begin{itemize}
    \item Provide all the project specifications for the proper implementation of the solving algorithm
	\item Provide the details to run an actual simulation of the system
\end{itemize}

\bibliographystyle{plainnat}
\bibliography{scoping_bib}

\end{document}
